\documentclass[11pt]{amsart}
\usepackage{amsmath,amssymb,amsthm,mathtools}
\usepackage[margin=1.15in]{geometry}
\usepackage{booktabs,longtable,array}
\usepackage{microtype}
\usepackage[hidelinks]{hyperref}
\usepackage{enumitem}
\setlist{leftmargin=2em}

\newcommand{\repo}{\url{https://github.com/jaideepsaipadhi/the-auluck-chowla-gupta-conjecture}}
\newcommand{\script}[1]{\texttt{#1}}

\newtheorem{theorem}{Theorem}[section]
\newtheorem{lemma}[theorem]{Lemma}
\newtheorem{proposition}[theorem]{Proposition}
\newtheorem{corollary}[theorem]{Corollary}
\newtheorem*{maintheorem}{Main Theorem}
\theoremstyle{definition}
\newtheorem{definition}[theorem]{Definition}
\theoremstyle{remark}
\newtheorem{remark}[theorem]{Remark}

\numberwithin{equation}{section}

\newcommand{\RR}{\mathbb{R}}
\newcommand{\CC}{\mathbb{C}}
\newcommand{\ZZ}{\mathbb{Z}}
\newcommand{\Li}{\operatorname{Li}}
\newcommand{\TV}{\operatorname{TV}}
\newcommand{\Var}{\operatorname{Var}}
\newcommand{\erfc}{\operatorname{erfc}}
\newcommand{\sinc}{\operatorname{sinc}}
\newcommand{\re}{\operatorname{Re}}
\newcommand{\im}{\operatorname{Im}}
\newcommand{\E}{\mathbb{E}}
\newcommand{\Prob}{\mathbb{P}}
\newcommand{\dd}{\,\mathrm{d}}
\newcommand{\ii}{\mathrm{i}}
\newcommand{\ee}{\mathrm{e}}

\title[The Auluck--Chowla--Gupta conjecture]{A proof of the Auluck--Chowla--Gupta Conjecture}
\author{Jaideep Sai Padhi}
\address{Purdue University, West Lafayette, IN 47907, USA}
\email{jpadhi@purdue.edu}
\subjclass[2020]{Primary 05A17; Secondary 11P82, 05A20, 65G30}
\keywords{integer partitions, unimodality, saddle-point method, interval arithmetic, computer-assisted proof}

\begin{document}

\begin{abstract}
Let $p(n,k)$ denote the number of partitions of $n$ into exactly $k$ parts. In 1942 Auluck, Chowla and Gupta
conjectured that for every $n$ the sequence $p(n,1),p(n,2),\dots,p(n,n)$ is unimodal. Szekeres proved this for all
sufficiently large $n$, but with a threshold that was never made explicit. We prove the conjecture for every $n\ge 1$.
For $n\ge 10^5$ we give an analytic argument with fully explicit constants: an exact probabilistic representation of the
ratio $p(n,k+1)/p(n,k)$, a second-order saddle-point main term with a rigorous uniform error bound in the transition window,
and separate elementary or saddle-point estimates in the two tails. Every numerical constant in this argument is certified in
ball arithmetic. For $n\le 2\cdot10^5$ the conjecture is verified exhaustively by two independent programs, one of them using
exact integer arithmetic only. Since log-concavity of $k\mapsto p(n,k)$ fails for every $n\ge 4$, the argument works with
first differences and never uses log-concavity.
\end{abstract}

\maketitle

%======================================================================
\section{Introduction}\label{sec:intro}
%======================================================================

For integers $n\ge1$ and $1\le k\le n$ let $p(n,k)$ be the number of partitions of $n$ into exactly $k$ positive parts. A finite
sequence $a_1,\dots,a_N$ of real numbers is \emph{(weakly) unimodal} if there is an index $t$ with
$a_1\le\dots\le a_t\ge a_{t+1}\ge\dots\ge a_N$.

\begin{maintheorem}
For every integer $n\ge1$ the sequence $p(n,1),p(n,2),\dots,p(n,n)$ is unimodal.
\end{maintheorem}

\subsection*{History}
The statement was conjectured by Auluck, Chowla and Gupta~\cite{ACG42}. Erd\H{o}s~\cite{Erdos46}, who studied the distribution of
the number of parts of a random partition (see also Erd\H{o}s--Lehner~\cite{EL41}), wrote that he could neither prove nor disprove
it. Szekeres~\cite{Sz51,Sz53} obtained an asymptotic formula for $p(n,k)$ uniform in $k$ and deduced the conjecture for all
$n$ larger than some $N_0$; the argument does not give a value of $N_0$ (see also~\cite{Sz90}). Canfield~\cite{Can97} later
gave a different derivation of Szekeres' formula, by recursions rather than by contour integration. To our knowledge no proof
for all $n$ has appeared before. % VERIFY: wording of the Szekeres/Canfield attribution against the original papers.

\subsection*{Why first differences}
Unimodality would follow from log-concavity, $p(n,k)^2\ge p(n,k-1)p(n,k+1)$, but this fails at the tail of every row with
$n\ge4$: for $n/2\le k\le n$ one has $p(n,k)=p(n-k)$ (Lemma~\ref{lem:E}), so $p(n,n)=p(n,n-1)=1$ and $p(n,n-2)=2$, and
$p(n,n-1)^2=1<2=p(n,n)p(n,n-2)$. We therefore work directly with the sign of the first difference, through the ratio
\[
  R_k=R_k(n):=\frac{p(n,k+1)}{p(n,k)}\qquad(1\le k\le n-1).
\]
The row is unimodal if and only if the signs of $\log R_k$, read left to right, never change from negative to positive.

\subsection*{Strategy}
Put
\begin{equation}\label{eq:beta}
  \beta:=\frac{\sqrt{6n}}{\pi},\qquad \ell:=\log\beta,\qquad \mu:=\beta\ell .
\end{equation}
The mode of the row lies within $O(\beta)$ of $\mu$ (this is the Gumbel scaling of Erd\H{o}s and Lehner~\cite{EL41}).
For $n\ge10^5$ we split $1\le k\le n-1$ into five ranges,
\[
\begin{gathered}
  D=[1,k_D],\qquad CL=(k_D,\mu-2\beta],\qquad W=(\mu-2\beta,\mu+2\beta),\\
  CR=[\mu+2\beta,\lceil n/2\rceil],\qquad E=[\lceil n/2\rceil,n-1],
\end{gathered}
\]
with $k_D=\lfloor1.7n^{1/3}\rfloor$, and determine the sign of $\log R_k$ on each range separately.

\begin{itemize}
\item On $D$ an elementary binomial sandwich shows $R_k>1$ (Section~\ref{sec:D}).
\item On $E$ an exact bijection shows $R_k\le1$ (Lemma~\ref{lem:E}).
\item Everything else rests on an exact identity (Section~\ref{sec:rep}): $R_k$ equals the ratio of two Fourier integrals
  of the characteristic function of a sum of independent geometric variables, taken at its saddle point $\theta$. This
  reduces the problem to estimating a ratio of integrals with the \emph{same} weight, so that all exponentially large factors
  cancel exactly.
\item On $CL$ and $CR$ a zeroth-order version of this estimate is enough to give $R_k>1$ and $R_k<1$ respectively
  (Section~\ref{sec:CLCR}).
\item In the window $W$, where the sign changes, we need more. We define an explicit second-order main term $L_{2,k}$
  (Section~\ref{sec:main}) and prove $|\log R_k-L_{2,k}|\le 0.4282\,\theta q$, where $q=e^{-(k+1)\theta}$ (Theorem~\ref{thm:A}).
  We then show that $L_2$ decreases by at least $0.91\,\theta q$ at each step (Theorem~\ref{thm:B}). Since
  $0.91>2\times0.4282$, at most one $k\in W$ has an undetermined sign, and that $k$ is adjacent to the sign change of $L_2$.
\end{itemize}

The constants in these estimates are bounded by rigorous ball arithmetic (Arb~\cite{Arb}, through python-flint) over finite
covers of the relevant parameter ranges. The analytic argument is carried out for $n\ge10^5$, and the range $n\le2\cdot10^5$
is checked exhaustively (Section~\ref{sec:finite}).

\subsection*{Computer assistance}
All computer-assisted steps are listed in Section~\ref{sec:computer} (Table~\ref{tab:certs}), each with the script that
certifies it and the exact statement certified. The scripts, their expected outputs and a driver that re-runs everything and
compares outputs are available at \repo. The mathematical arguments in this paper do not depend on any computation other than
those listed there.

\subsection*{Organisation}
Section~\ref{sec:rep} gives the notation and the exact representation. Section~\ref{sec:deriv} proves uniform derivative bounds.
Section~\ref{sec:main} defines the main term and states and proves the error bound of Theorem~\ref{thm:A}; its uniformity rests on
the ingredient bounds of Section~\ref{sec:ingredients}. Section~\ref{sec:B} analyses the step of $L_2$. Section~\ref{sec:CLCR}
treats the tail regions, Section~\ref{sec:assembly} assembles the proof for $n\ge10^5$, Section~\ref{sec:finite} describes the
finite verification, and Section~\ref{sec:computer} lists the computer-assisted inequalities and the trust base. The longer
error algebra is in the appendices.

\subsection*{Conventions}
$\zeta_2:=\zeta(2)=\pi^2/6$, so that $n=\zeta_2\beta^2$. $\Li_s(y)=\sum_{l\ge1}y^l/l^s$ is the polylogarithm; for integers
$r\ge1$, $\Li_{1-r}(y)=\sum_l l^{r-1}y^l$ is a rational function of $y$. Logarithms of complex numbers are principal and are only
taken of numbers with positive real part. For a function $g$ and an interval $I$, $\Var_I g$ is the total variation of $g$ on $I$,
and $\TV(g)=\Var_{[0,\infty)}g$.

%======================================================================
\section{Notation and an exact reduction}\label{sec:rep}
%======================================================================

Fix $n$ and $2\le k\le n-2$. Write $a_k=p(n,k)$ and
\[
  m:=n-k\ \ (\ge2),\qquad K:=k+1 .
\]
Let $f_k(N)$ be the number of partitions of $N$ into parts of size at most $k$, with $f_k(N)=0$ for $N<0$.

\begin{lemma}\label{lem:comb}
$p(n,k)=f_k(n-k)$, and $f_{k+1}(N)=\sum_{s\ge0}f_k(N-sK)$. Consequently
$a_{k+1}=\sum_{s\ge0}f_k(m-1-sK)$, a finite sum.
\end{lemma}

\begin{proof}
These are standard (see~\cite{Andrews}). The conjugate of a partition of $n$ into exactly $k$ parts has largest part exactly $k$; removing one part $k$ is a bijection
onto partitions of $n-k$ into parts $\le k$. For the second identity, classify partitions of $N$ into parts $\le K$ by the
number $s$ of parts equal to $K$ and remove them. Finally $a_{k+1}=f_{k+1}(m-1)$.
\end{proof}

\subsection{The saddle point and a random variable}
For $t>0$ let
\begin{equation}\label{eq:S}
  S_k(t):=\sum_{i=1}^{k}\frac{i}{e^{it}-1}.
\end{equation}
$S_k$ is continuous and strictly decreasing from $+\infty$ to $0$, so there is a unique $\theta=\theta(n,k)>0$ with
$S_k(\theta)=m$. We call $\theta$ the \emph{saddle point} and put
\begin{equation}\label{eq:params}
  x:=e^{-\theta},\quad q:=x^K=e^{-K\theta},\quad \nu:=K\theta,\quad \nu':=\nu+\log\theta,\quad v:=k\theta=\nu-\theta,
\end{equation}
so that $q=\theta e^{-\nu'}$. Let $Y_1,\dots,Y_k$ be independent with $\Prob(Y_i=y)=(1-x^i)x^{iy}$ $(y\ge0)$ and put
$X:=\sum_{i\le k}iY_i$.

\begin{lemma}\label{lem:X}
$\Prob(X=N)=f_k(N)x^N\prod_{i\le k}(1-x^i)$, and $\E X=S_k(\theta)=m$.
\end{lemma}

\begin{proof}
$\Prob(X=N)=\prod_i(1-x^i)\cdot x^N\cdot\#\{(y_i):\sum iy_i=N\}$ and the count is $f_k(N)$, with $y_i$ the multiplicity of the
part $i$. Also $\E Y_i=x^i/(1-x^i)=1/(e^{i\theta}-1)$.
\end{proof}

\begin{proposition}[Exact representation]\label{prop:rep}
Let $\chi(\varphi):=\E\,e^{\ii\varphi(X-m)}$ and
\[
  g_0:=\frac{x}{1-q},\qquad H(\varphi):=\ii\varphi-\log\bigl(1-qe^{\ii K\varphi}\bigr)+\log(1-q).
\]
Then
\begin{equation}\label{eq:rep1}
  R_k=\sum_{s\ge0}x^{1+sK}\,\frac{\Prob(X=m-1-sK)}{\Prob(X=m)},
\end{equation}
and, with $D:=\int_{-\pi}^{\pi}\chi(\varphi)\dd\varphi=2\pi\Prob(X=m)>0$ and $\tilde N:=\int_{-\pi}^{\pi}\chi(\varphi)\bigl(e^{H(\varphi)}-1\bigr)\dd\varphi$,
\begin{equation}\label{eq:rep2}
  \frac{R_k}{g_0}-1=\frac{\tilde N}{D}.
\end{equation}
\end{proposition}

\begin{proof}
Since $f_k(m)\ge1$ we have $\Prob(X=m)>0$. By Lemma~\ref{lem:X}, $f_k(N)=\Prob(X=N)x^{-N}/\prod(1-x^i)$, and inserting this into
Lemma~\ref{lem:comb} gives~\eqref{eq:rep1}. The series $\chi(\varphi)=\sum_N\Prob(X=N)e^{\ii\varphi(N-m)}$ converges absolutely and
uniformly, so $\frac1{2\pi}\int_{-\pi}^{\pi}\chi(\varphi)e^{\ii j\varphi}\dd\varphi=\Prob(X=m-j)$. Inserting this into~\eqref{eq:rep1} and
exchanging the sum over $s$ with the integral (dominated by $\sum_s x^{1+sK}$) gives
$R_k=\int\chi G/\int\chi$ with $G(\varphi)=\sum_s x^{1+sK}e^{\ii(1+sK)\varphi}=xe^{\ii\varphi}/(1-qe^{\ii K\varphi})$. Since
$\re(1-qe^{\ii K\varphi})\ge1-q>0$, $H$ is well defined and $e^H=G/g_0$, which gives~\eqref{eq:rep2}.
\end{proof}

\begin{remark}
Identity~\eqref{eq:rep1} holds for every $x\in(0,1)$, not only at the saddle point, because the $x$-dependence cancels. The
saddle point is chosen so that $\chi$ is concentrated near $\varphi=0$. The form~\eqref{eq:rep2}, rather than $R_k=N/D$, matters:
a relative error $\delta$ in $D$ costs only $\delta\,|R_k/g_0-1|$, which is of order $\delta\theta^2$.
\end{remark}

By independence,
\begin{equation}\label{eq:psi}
  \chi=e^{\psi},\qquad \psi(\varphi):=-\ii m\varphi+\sum_{i\le k}\Bigl[\log(1-x^i)-\log\bigl(1-x^ie^{\ii i\varphi}\bigr)\Bigr],
\end{equation}
a real-analytic function of $\varphi\in\RR$. Since $X$ is real-valued, $\chi(-\varphi)=\overline{\chi(\varphi)}$, so
$\re\psi$ is even and $\im\psi$ is odd. Moreover
\begin{equation}\label{eq:modchi}
  -2\log|\chi(\varphi)|=S(\varphi):=\sum_{i\le k}\log\Bigl(1+\frac{\sin^2(i\varphi/2)}{\sinh^2(i\theta/2)}\Bigr),
\end{equation}
because $|1-ye^{\ii a}|^2/(1-y)^2=1+\sin^2(a/2)/\sinh^2(u/2)$ for $y=e^{-u}$. Every summand of $S$ is nonnegative.

\subsection{Cumulants and the functions \texorpdfstring{$f_r$, $F_r$}{f\_r, F\_r}}
For $r\ge2$ put
\begin{equation}\label{eq:kappa}
  \kappa_r:=\sum_{i\le k}i^r\Li_{1-r}(x^i),\qquad B:=\kappa_2,\qquad
  c_1:=1+\frac{Kq}{1-q},\qquad c_r:=K^r\Li_{1-r}(q)\ (r\ge2).
\end{equation}
For $r\ge1$ let $f_r(u):=u^r\Li_{1-r}(e^{-u})$ for $u>0$, extended continuously by $f_r(0)=(r-1)!$, and let
$F_r(v):=\int_0^vf_r$. Thus $f_1(u)=u/(e^u-1)$ and $f_2(u)=(u/2)^2/\sinh^2(u/2)$; both are decreasing and take values in $(0,1]$.
We use the scaled quantities
\begin{equation}\label{eq:scaled}
  s_r:=\theta^{r+1}\kappa_r=\theta\sum_{i\le k}f_r(i\theta),\qquad b:=s_2=\theta^3B,\qquad c_r\theta^r=f_r(\nu)\ (r\ge2),
  \qquad c_1\theta=\theta+f_1(\nu).
\end{equation}

\begin{lemma}[Closed form of $F_r$]\label{lem:Fr}
For $r\ge1$ and $v>0$,
\[
  F_r(v)=r!\Bigl(\zeta_2-\sum_{j=0}^{r}\frac{v^j}{j!}\Li_{2-j}(e^{-v})\Bigr),\qquad F_r(\infty)=r!\,\zeta_2 .
\]
In particular $F_r$ is nondecreasing and bounded by $r!\zeta_2$.
\end{lemma}

\begin{proof}
Call the right side $G_r$. From $\frac{d}{dv}\Li_s(e^{-v})=-\Li_{s-1}(e^{-v})$ and the Leibniz rule, the derivative of the sum over $j$
telescopes to $-v^r\Li_{1-r}(e^{-v})/r!$, so $G_r'=f_r$. As $v\to0^+$, $\Li_2(e^{-v})\to\zeta_2$, $v\Li_1(e^{-v})=-v\log(1-e^{-v})\to0$,
and $v^j\Li_{2-j}(e^{-v})=O(v)$ for $j\ge2$ because $\Li_{2-j}(e^{-v})=O(v^{1-j})$; hence $G_r(0^+)=0=F_r(0)$. Since $f_r$ is continuous
on $[0,v]$, $F_r=G_r$. As $v\to\infty$ every $v^j\Li_{2-j}(e^{-v})$ tends to $0$.
\end{proof}

\subsection{Main term}
Define
\begin{equation}\label{eq:M2}
\begin{aligned}
  N_0&:=1+\frac{\kappa_4}{8B^2}-\frac{5\kappa_3^2}{24B^3},\qquad
  E_2:=\frac1B\Bigl(1+\frac{\kappa_4}{2B^2}-\frac{5\kappa_3^2}{4B^3}\Bigr),\\
  E_1&:=\frac1{N_0}\Bigl(\frac{\kappa_3}{2B^2}-\frac{\kappa_5}{8B^3}+\frac{35\,\kappa_3\kappa_4}{48\,B^4}-\frac{35\,\kappa_3^3}{48\,B^5}\Bigr),\\
  M_2&:=1+c_1E_1-\frac{(c_2+c_1^2)E_2}{2}-\frac{5c_3\kappa_3}{12B^3}+\frac{c_4}{8B^2},
\end{aligned}
\end{equation}
and
\begin{equation}\label{eq:L2}
  L_k:=\log R_k,\qquad L_{2,k}:=-\theta-\log(1-q)+\log M_2 = \log g_0+\log M_2 .
\end{equation}
The origin of $M_2$ is explained in Section~\ref{sec:main}: it is the closed form of the ratio of the Gaussian model integrals
up to the order needed. Note that $L_k-L_{2,k}=\log\bigl(R_k/(g_0M_2)\bigr)$.

%======================================================================
\section{Uniform derivative bounds}\label{sec:deriv}
%======================================================================

\begin{lemma}\label{lem:psider}
$\psi(0)=\psi'(0)=0$, $\psi^{(r)}(0)=\ii^r\kappa_r$ and $|\psi^{(r)}(\varphi)|\le\kappa_r$ for all $\varphi\in\RR$ and $r\ge2$.
\end{lemma}

\begin{proof}
For $0<y<1$ and $a>0$, $\ell_{y,a}(\varphi):=-\log(1-ye^{\ii a\varphi})=\sum_{l\ge1}y^le^{\ii la\varphi}/l$, absolutely and uniformly.
The $r$-times differentiated series $(\ii a)^r\sum_ll^{r-1}y^le^{\ii la\varphi}$ is dominated by $a^r\Li_{1-r}(y)<\infty$, so by the
Weierstrass M-test termwise differentiation is valid and $|\ell_{y,a}^{(r)}|\le a^r\Li_{1-r}(y)$. Now
$\psi=-\ii m\varphi+\sum_i[\log(1-x^i)+\ell_{x^i,i}]$; for $r\ge2$ the linear part disappears, which gives the bound and the value at $0$.
Finally $\psi'(0)=\ii\bigl(-m+\sum_i ix^i/(1-x^i)\bigr)=0$ by the saddle-point equation.
\end{proof}

\begin{lemma}\label{lem:Hder}
$H(0)=0$, $H'(0)=\ii c_1$, $H^{(r)}(0)=\ii^rc_r$ for $r\ge2$; $|H'|\le c_1$ and $|H^{(r)}|\le c_r$ $(r\ge2)$ on $\RR$; and $\re H\le0$.
\end{lemma}

\begin{proof}
$H=\ii\varphi+\ell_{q,K}+\log(1-q)$, so $H'=\ii+\ii K\sum_lq^le^{\ii lK\varphi}$ and $|H'|\le1+Kq/(1-q)=c_1$; for $r\ge2$,
$H^{(r)}=\ell_{q,K}^{(r)}$. Finally $|1-qe^{\ii K\varphi}|^2-(1-q)^2=2q(1-\cos K\varphi)\ge0$, so $|G|\le g_0$ and $\re H=\log|G/g_0|\le0$.
\end{proof}

\begin{corollary}\label{cor:omega}
Let $\omega:=\psi+B\varphi^2/2$. Then $\omega(0)=\omega'(0)=\omega''(0)=0$ and $|\omega'''|\le\kappa_3$; hence $|\omega(\varphi)|\le\kappa_3|\varphi|^3/6$.
Moreover $|\im\psi(\varphi)|\le\kappa_3|\varphi|^3/6$ and $\re\psi(\varphi)\le-B\varphi^2/2+\kappa_4\varphi^4/24$.
\end{corollary}

\begin{proof}
The first statements follow from Lemma~\ref{lem:psider} and Taylor's theorem with integral remainder (applied to real and
imaginary parts). $\im\psi$ is odd and vanishes to second order at $0$ because $\psi'(0)=0$ and $\psi''(0)=-B$ is real.
$\re\psi$ is even, with $(\re\psi)''(0)=-B$, vanishing odd derivatives at $0$, and $|(\re\psi)''''|\le\kappa_4$.
\end{proof}

We shall also use two elementary facts repeatedly.
\begin{enumerate}[label=(F\arabic*)]
\item\label{F1} If $\re z\le0$ then $|e^z-1|\le|z|$ and $|e^z-1-z-z^2/2|\le|z|^3/6$; more generally
  $|e^z-1|\le|z|\max(1,e^{\re z})$. (From $e^z-1=z\int_0^1e^{sz}\dd s$ and $e^z-1-z-z^2/2=z^3\int_0^1\frac{(1-s)^2}{2}e^{sz}\dd s$.)
\item\label{F2} For $0\le A$ and $0\le t\le1$, $\log(1+At)\ge t\log(1+A)$ (concavity).
\end{enumerate}

%======================================================================
\section{The main term and the error bound in the window}\label{sec:main}
%======================================================================

Throughout this section $2\le k\le n-2$ and $\theta$ is the saddle point. We fix the central arc
\[
  |\varphi|\le\varphi_0:=0.4\,\theta .
\]

\subsection{The Gaussian model}
Put $g(\varphi):=e^{-B\varphi^2/2}$, and let
\[
\begin{gathered}
  \omega_3=-\ii\kappa_3\varphi^3/6,\qquad \omega_4=\kappa_4\varphi^4/24,\qquad \omega_5=\ii\kappa_5\varphi^5/120,\\
  H_1=\ii c_1\varphi,\qquad H_2=-c_2\varphi^2/2,\qquad H_3=-\ii c_3\varphi^3/6,\qquad H_4=c_4\varphi^4/24,
\end{gathered}
\]
be the Taylor terms of $\omega=\psi+B\varphi^2/2$ and of $H$ at $0$ (Lemmas~\ref{lem:psider} and~\ref{lem:Hder}). Define the
polynomials
\[
  Q:=1+\omega_3+\omega_4+\omega_5+\tfrac12\omega_3^2+\omega_3\omega_4+\tfrac16\omega_3^3,\qquad
  P_H:=1+H_1+H_2+H_3+H_4+\tfrac12H_1^2,
\]
and the model integrals
\[
  D_0:=\int_\RR g\,Q\dd\varphi,\qquad \tilde N_0:=\int_\RR g\,Q\,(P_H-1)\dd\varphi,\qquad M_{\mathrm{model}}:=1+\tilde N_0/D_0 .
\]
They are evaluated exactly with the Gaussian moments $\int_\RR\varphi^{2j}e^{-A\varphi^2/2}\dd\varphi=\sqrt{2\pi/A}\,(2j-1)!!\,A^{-j}$.
In particular $D_0=\sqrt{2\pi/B}\,N_0$ with $N_0$ as in~\eqref{eq:M2}, and $\tilde N_0/D_0$ is real. We work with the scaled variables
\begin{equation}\label{eq:AC}
  A_r:=\frac{\kappa_r}{B^{r/2}}=\frac{s_r\,\theta^{r/2-1}}{b^{r/2}},\qquad C_r:=\frac{c_r}{B^{r/2}}=\frac{f_r(\nu)\,\theta^{r/2}}{b^{r/2}}\ \ (r\ge2),\qquad
  C_1:=\frac{c_1}{B^{1/2}},
\end{equation}
and $w_0:=\varphi_0\sqrt B=0.4\sqrt{b/\theta}$; every quantity below is invariant under $\kappa_r\mapsto\lambda^r\kappa_r$,
$c_r\mapsto\lambda^rc_r$, $B\mapsto\lambda^2B$, so it may be computed with $B=1$ and $(\kappa_r,c_r)$ replaced by $(A_r,C_r)$.

\begin{lemma}\label{lem:gap}
$M_{\mathrm{model}}-M_2=\mathrm{Num}/N_0$, where, in the scaled variables,
\[
\begin{aligned}
\mathrm{Num}={}&-\tfrac{25}{48}C_4A_3^2+\tfrac16C_4A_4+(C_2+C_1^2)\Bigl(\tfrac{1}{32}A_4^2+\tfrac{25}{192}A_3^4-\tfrac{25}{192}A_3^2A_4\Bigr)\\
&+\tfrac54C_3A_3^3+\tfrac{7}{48}C_3A_5-\tfrac{25}{24}C_3A_3A_4 .
\end{aligned}
\]
\end{lemma}

\begin{proof}
A polynomial identity in the variables $A_3,\dots,A_5,C_1,\dots,C_4$, obtained by expanding $\tilde N_0$ and $D_0$ in Gaussian
moments. It is verified symbolically (Table~\ref{tab:certs}, item~S2).
\end{proof}

On the window, $f_r(\nu)=O(q)=O(\theta)$, so by~\eqref{eq:AC} $A_r=O(\theta^{r/2-1})$, $C_r=O(\theta^{1+r/2})$ $(r\ge2)$ and
$C_1=O(\theta^{3/2})$. Hence every monomial of $\mathrm{Num}$ is $O(\theta^4)$, whereas $M_2-1=O(\theta^2)$ is of the same order as the
step $\theta q$ of $L_2$ (Theorem~\ref{thm:B}); this is why a second-order main term is needed.

\subsection{The pointwise error bound}
Let $a_r:=\kappa_r/r!$ and $t=|\varphi|$, and set
\[
  \delta_0:=\frac{\kappa_4\varphi_0^2/12+\kappa_6\varphi_0^4/360}{B},\qquad B':=(1-\delta_0)B .
\]
Define the nonnegative polynomials in $t$
\[
\begin{aligned}
T_\omega(t)&:=a_6t^6+\bigl(\tfrac12a_4^2+a_3a_5\bigr)t^8+a_4a_5t^9+\tfrac12a_5^2t^{10}\\
  &\qquad+\tfrac12(a_4t^4+a_5t^5)(a_3t^3+a_4t^4+a_5t^5)^2+\tfrac1{24}(a_3t^3+a_4t^4+a_5t^5)^4,\\
T_H(t)&:=\tfrac{1}{120}c_5t^5+\bigl(\tfrac12c_1c_2+\tfrac16c_1^3\bigr)t^3,
\end{aligned}
\]
and let $Q_{\mathrm{abs}}(t)$ and $P_{\mathrm{abs}}(t)$ be the polynomials obtained from $Q$ and $P_H-1$ by replacing each coefficient by
its modulus.

\begin{lemma}[Central arc]\label{lem:central}
Suppose $\delta_0<1$. For $|\varphi|\le\varphi_0$,
\[
  |e^{\omega}-Q|\le T_\omega(t)\,e^{\delta_0Bt^2/2},\qquad |e^H-P_H|\le T_H(t),\qquad |e^H-1|\le c_1t,
\]
and consequently
\[
\begin{aligned}
  \Bigl|\int_{|\varphi|\le\varphi_0}\bigl[\chi(e^H-1)-gQ(P_H-1)\bigr]\Bigr|&\le\int_\RR e^{-B't^2/2}c_1t\,T_\omega(t)\dd\varphi+\int_\RR g\,Q_{\mathrm{abs}}T_H\dd\varphi,\\
  \Bigl|\int_{|\varphi|\le\varphi_0}(\chi-gQ)\Bigr|&\le\int_\RR e^{-B't^2/2}\,T_\omega(t)\dd\varphi .
\end{aligned}
\]
\end{lemma}

\begin{proof}
By Taylor's theorem with $|\psi^{(6)}|\le\kappa_6$ and since the odd-order terms are purely imaginary,
$\re\omega\le\kappa_4t^4/24+\kappa_6t^6/720\le\delta_0Bt^2/2$ for $t\le\varphi_0$. Let $\hat\omega=\omega_3+\omega_4+\omega_5$; then
$|\omega-\hat\omega|\le a_6t^6$ and $\re\hat\omega=\kappa_4t^4/24\le\delta_0Bt^2/2$. We split
\[
  e^\omega-Q=(e^\omega-e^{\hat\omega})+\Bigl(e^{\hat\omega}-\sum_{j\le3}\tfrac{\hat\omega^j}{j!}\Bigr)+\Bigl(\sum_{j\le3}\tfrac{\hat\omega^j}{j!}-Q\Bigr).
\]
The first term is at most $|\omega-\hat\omega|e^{\max(\re\omega,\re\hat\omega)}$, the second is at most
$\frac{|\hat\omega|^4}{24}\max(1,e^{\re\hat\omega})$ (Taylor's theorem with integral remainder), and the third is the polynomial
$\frac12(\omega_4^2+\omega_5^2)+\omega_3\omega_5+\omega_4\omega_5+\frac16(\hat\omega^3-\omega_3^3)$, whose modulus is dominated termwise by
the remaining terms of $T_\omega$. For $T_H$, write
$e^H-P_H=(e^H-1-H-\frac12H^2)+(H-\sum_{r\le4}H_r)+\frac12(H-H_1)(H+H_1)$. By Lemma~\ref{lem:Hder}, $|H|\le c_1t$, so
\ref{F1} bounds the first piece by $c_1^3t^3/6$; Taylor's theorem with $|H^{(5)}|\le c_5$ bounds the second by $c_5t^5/120$; and
$|H-H_1|\le c_2t^2/2$, $|H+H_1|\le2c_1t$ bound the third by $c_1c_2t^3/2$. The bound $|e^H-1|\le c_1t$ is~\ref{F1}. The integral
inequalities follow from the identity
$\chi(e^H-1)-gQ(P_H-1)=g\bigl[(e^\omega-Q)(e^H-1)+Q(e^H-P_H)\bigr]$ and from $g\,e^{\delta_0Bt^2/2}=e^{-B't^2/2}$.
\end{proof}

The \emph{model tails} are bounded by $e^{-B\varphi^2/2}\le e^{-B\varphi_0^2/4}e^{-B\varphi^2/4}$ for $|\varphi|\ge\varphi_0$:
\begin{equation}\label{eq:tails}
\begin{aligned}
  \int_{|\varphi|>\varphi_0}g\,|Q(P_H-1)|&\le e^{-B\varphi_0^2/4}\int_\RR e^{-B\varphi^2/4}Q_{\mathrm{abs}}P_{\mathrm{abs}},\\
  \int_{|\varphi|>\varphi_0}g\,|Q|&\le e^{-B\varphi_0^2/4}\int_\RR e^{-B\varphi^2/4}Q_{\mathrm{abs}} .
\end{aligned}
\end{equation}
On the \emph{minor arcs} $\varphi_0\le|\varphi|\le\pi$, since $|e^H-1|\le2$ by $\re H\le0$ and the arcs have total length at most $2\pi$,
\begin{equation}\label{eq:minor}
  \Bigl|\int_{\varphi_0\le|\varphi|\le\pi}\chi(e^H-1)\Bigr|\le4\pi M_\chi,\qquad \Bigl|\int_{\varphi_0\le|\varphi|\le\pi}\chi\Bigr|\le2\pi M_\chi,
\end{equation}
where $M_\chi:=\sup_{\varphi_0\le|\varphi|\le\pi}|\chi(\varphi)|$.

\begin{proposition}[Pointwise error bound]\label{prop:eps}
Let $E_N$ (resp.\ $E_D$) be the sum of the numerator (resp.\ denominator) bounds of Lemma~\ref{lem:central},~\eqref{eq:tails}
and~\eqref{eq:minor}, all divided by $\sqrt{2\pi/B}$ (recall $D_0/\sqrt{2\pi/B}=N_0$). Suppose $\delta_0<1$, $N_0-E_D>0$ and
\[
  \Delta:=\frac{E_N+|\tilde N_0/D_0|\,E_D}{N_0-E_D}+\frac{|\mathrm{Num}|}{N_0}<M_2 .
\]
Then
\begin{equation}\label{eq:eps}
  |L_k-L_{2,k}|\le\varepsilon_k:=-\log\bigl(1-\Delta/M_2\bigr).
\end{equation}
\end{proposition}

\begin{proof}
Write $\tilde N=\tilde N_0+e_N$ and $D=D_0+e_D$. The three groups of estimates give $|e_N|\le\sqrt{2\pi/B}\,E_N$ and
$|e_D|\le\sqrt{2\pi/B}\,E_D$. Since
$\tilde N/D-\tilde N_0/D_0=(e_N-(\tilde N_0/D_0)e_D)/(D_0+e_D)$, Proposition~\ref{prop:rep} gives
$|R_k/g_0-M_{\mathrm{model}}|\le\Delta-|\mathrm{Num}|/N_0$, and Lemma~\ref{lem:gap} gives $|R_k/(g_0M_2)-1|\le\Delta/M_2<1$. Finally
$|\log(1+y)|\le-\log(1-|y|)$ for $|y|<1$, and $L_k-L_{2,k}=\log(R_k/(g_0M_2))$.
\end{proof}

All the integrals in $E_N$ and $E_D$ are integrals of nonnegative polynomials against Gaussians over $\RR$, which we evaluate exactly
via $\int_\RR|\varphi|^pe^{-A\varphi^2/2}\dd\varphi=(2/A)^{(p+1)/2}\Gamma(\frac{p+1}2)$. The full expression for $\varepsilon_k$, written
in the scaled variables~\eqref{eq:AC}, is recorded in Appendix~\ref{app:eps}. It is a fixed algebraic expression in
$(\theta,\nu',b,s_3,\dots,s_6)$ and in an upper bound for $M_\chi$; all of its algebra (the moments, the polynomial
identities above, and the scaling invariance) is checked symbolically (Table~\ref{tab:certs}, S2).

\subsection{The minor arcs}
The only non-local input is a bound for $M_\chi$. The following lemma gives it uniformly in $\theta$; its proof uses only the
partial sums of $\sin^2$, and no counting arguments.

\begin{lemma}[Minor arcs]\label{lem:minor}
Let $\theta_{\max}>0$, $u_{\mathrm{top}}>0$ and $c_0>0$, and suppose $\theta\le\theta_{\max}$ and $(k+1)\theta\ge u_{\mathrm{top}}$.
Let $c_a\ge c_0$ with $c_a\theta_{\max}/2\le\pi/2$, and put $L_\ast:=2\theta_{\max}+2/\bigl(c_a\sinc(c_a\theta_{\max}/2)\bigr)$.
\begin{enumerate}[label=(\alph*)]
\item If $\varphi/\theta\in[c_{\mathrm{lo}},c_{\mathrm{hi}}]\subset[c_0,c_a]$, then for every $T\in(0,\pi/2]$, with $\sigma=\sin T/T$,
\[
  \theta S(\varphi)\ \ge\ \int_{\theta_{\max}}^{\min(u_{\mathrm{top}},\,2T/c_{\mathrm{hi}})}\log\bigl(1+\sigma^2c_{\mathrm{lo}}^2f_2(u)\bigr)\dd u .
\]
\item If $c_a\theta\le\varphi\le\pi$, then $\theta S(\varphi)\ge\int_{L_\ast}^{u_{\mathrm{top}}}\log\coth(u/2)\dd u$ (to be read as $0$ if $L_\ast\ge u_{\mathrm{top}}$).
\end{enumerate}
Consequently $\log|\chi(\varphi)|\le-\Psi/(2\theta)$ for $c_0\theta\le|\varphi|\le\pi$, where $\Psi$ is the minimum of the right-hand
sides over a cover of $[c_0,c_a]$ by cells and of (b).
\end{lemma}

\begin{proof}
By~\eqref{eq:modchi} and $\chi(-\varphi)=\overline{\chi(\varphi)}$ we may take $\varphi>0$; recall that every summand of $S$ is $\ge0$.

(a) Put $c=\varphi/\theta$. For $i\le2T/\varphi$ we have $i\varphi/2\le T$, so $\sin(i\varphi/2)\ge\sigma\,i\varphi/2$, because $\sin y/y$ is
decreasing on $(0,\pi)$. Hence the $i$-th summand is at least
$\log\bigl(1+\sigma^2c^2(i\theta/2)^2/\sinh^2(i\theta/2)\bigr)\ge\lambda(i\theta)$ with $\lambda(u):=\log(1+\sigma^2c_{\mathrm{lo}}^2f_2(u))$,
a decreasing function. Keep only $i\le I:=\min(k,\lfloor2T/\varphi\rfloor)$. Then
$\theta S\ge\theta\sum_{i\le I}\lambda(i\theta)\ge\int_\theta^{(I+1)\theta}\lambda$. Since $(\lfloor2T/\varphi\rfloor+1)\theta>2T/c$, we have
$(I+1)\theta\ge\min((k+1)\theta,2T/c)\ge\min(u_{\mathrm{top}},2T/c_{\mathrm{hi}})$, and $\theta\le\theta_{\max}$, $\lambda\ge0$.

(b) Put $d=\varphi/2\in(0,\pi/2]$ and $w_i:=2\log\coth(i\theta/2)=\log(1+1/\sinh^2(i\theta/2))$, which is decreasing in $i$.
By~\ref{F2} with $t=\sin^2(id)$, the $i$-th summand is at least $\sin^2(id)\,w_i$. Let $G_I:=\sum_{i\le I}\sin^2(id)$. The Dirichlet
identity
\begin{equation}\label{eq:dirichlet}
  G_I=\frac I2-\frac{\sin(Id)\cos((I+1)d)}{2\sin d}
\end{equation}
gives $G_I\ge I/2-1/(2\sin d)$, and since also $G_I\ge0$ we get $G_I\ge\frac12(I-N_d)_+$ with $N_d:=\lceil1/\sin d\rceil$.
By Abel summation, with $w_{k+1}:=0$ and $w_I-w_{I+1}\ge0$,
\[
  S\ \ge\ \sum_{i\le k}\sin^2(id)w_i=\sum_{I\le k}G_I(w_I-w_{I+1})\ \ge\ \tfrac12\sum_{I\le k}(I-N_d)_+(w_I-w_{I+1})=\tfrac12\sum_{N_d<i\le k}w_i .
\]
As $w$ is decreasing, $w_i\ge\theta^{-1}\int_{i\theta}^{(i+1)\theta}2\log\coth(u/2)\dd u$, so
$\theta S\ge\int_{(N_d+1)\theta}^{(k+1)\theta}\log\coth(u/2)\dd u$. Finally $(N_d+1)\theta\le2\theta+\theta/\sin d$, and since $\sin$ is
increasing on $[0,\pi/2]$ and $\sinc$ is decreasing, $\sin d\ge\sin(c_a\theta/2)\ge(c_a\theta/2)\sinc(c_a\theta_{\max}/2)$. Hence
$(N_d+1)\theta\le L_\ast$ and $(k+1)\theta\ge u_{\mathrm{top}}$.
\end{proof}

\begin{corollary}\label{cor:cstar}
If $\theta\le0.0042$ and $(k+1)\theta\ge3$, then $|\chi(\varphi)|\le e^{-c_\ast/\theta}$ for $0.4\theta\le|\varphi|\le\pi$, with
$c_\ast=0.15973$.
\end{corollary}

\begin{proof}
Apply Lemma~\ref{lem:minor} with $\theta_{\max}=0.0042$, $u_{\mathrm{top}}=3$, $c_0=0.4$ and $c_a=0.4\cdot150^{96/400}=1.33144\ldots$;
then $L_\ast\le1.5105$. All integrals are right Riemann sums of decreasing integrands, hence rigorous lower bounds, evaluated in
ball arithmetic: (a) gives $\theta S\ge0.31946$ on the $96$ geometric cells of ratio $150^{1/400}$ covering $[0.4,c_a]$ (best of
$41$ values of $T$ per cell), and (b) gives
$\theta S\ge0.34433$. So $\Psi\ge0.31946$ and $c_\ast=\Psi/2\ge0.15973$ (Table~\ref{tab:certs}, C3).
\end{proof}

\subsection{Statement of Theorem A}
\begin{theorem}[Theorem A]\label{thm:A}
Let $n\ge10^5$ and $k\in W=(\mu-2\beta,\mu+2\beta)$. Then
\[
  |L_k-L_{2,k}|\le\varepsilon_k\le\Gamma_{\max}\,\theta q,\qquad \Gamma_{\max}=0.4282 .
\]
More precisely $\varepsilon_k\le0.42024\,\theta q$ when $\theta\ge0.002$, and $\varepsilon_k\le0.17605\,\theta q$ when $\theta\le0.002$.
\end{theorem}

The first inequality is Proposition~\ref{prop:eps}. The second is proved in Section~\ref{sec:proofA} after the ingredient bounds are in
place.

%======================================================================
\section{Ingredient bounds}\label{sec:ingredients}
%======================================================================

\subsection{The window}
Write $n=\zeta_2\beta^2$ and $\beta_A:=\sqrt{6\cdot10^5}/\pi=246.56\ldots$, $\ell_A=\log\beta_A$.

\begin{lemma}[Lemma W]\label{lem:W}
Let $n\ge10^5$, $k\in W$, and put $\delta:=5\ell/\beta$, $\epsilon_\beta:=(\ell+2+1/\beta)/(\zeta_2\beta)$ and $r:=(1-\epsilon_\beta)^{-1/2}$. Then
\[
  \frac{1-\delta}{\beta}<\theta\le\frac r\beta,
\]
and consequently
\[
  \theta\le\theta_A:=0.0040939,\qquad \nu'\in[-2.5102,\ 2.0840],\qquad \nu\ge3.1158 .
\]
\end{lemma}

\begin{proof}
\emph{Two-sided bound for $t^2S_k(t)$.} We have $t^2S_k(t)=t\sum_{i\le k}f_1(it)$ with $f_1$ decreasing and $0<f_1\le1$. Comparing with
integrals, $F_1((k+1)t)-F_1(t)\le t^2S_k(t)\le F_1(kt)\le\zeta_2$, and $F_1(t)\le t$. Write $F_1(v)=\zeta_2-T(v)$ with
$T(v)=\int_v^\infty f_1=\sum_l e^{-lv}(v/l+1/l^2)\le(v+1)e^{-v}/(1-e^{-v})$.

\emph{Upper bound.} Since $k<\mu+2\beta$, $m>n-\beta(\ell+2)$, and $\theta^2m\le\zeta_2$ gives
$\theta^2\le\zeta_2/m\le r^2/\beta^2$.

\emph{Lower bound.} Let $\theta_-:=(1-\delta)/\beta$. Since $S_k$ is decreasing it suffices to show $S_k(\theta_-)>m$. Now
$\theta_-^2m\le\theta_-^2n=(1-\delta)^2\zeta_2$ and $\theta_-^2S_k(\theta_-)\ge\zeta_2-T(K\theta_-)-\theta_-$, so it suffices that
$(2\delta-\delta^2)\zeta_2>T(K\theta_-)+\theta_-$. As $K>\beta(\ell-2)$, $K\theta_-\ge(\ell-2)(1-\delta)=:v_-$, and the bound for $T$
turns the sufficient condition into $G>0$, where (after multiplication by $\beta/\ell$)
\[
  G:=2\cdot5\,\zeta_2-25\,\zeta_2\,\ell/\beta-e^2\beta^{\delta}/(1-e^{-v_-})-1/\ell .
\]

\emph{Parameter ranges.} At fixed $K$, $\nu'=K\theta+\log\theta$ is increasing in $\theta$. With $K<\mu+2\beta+1$ and $\theta\le r/\beta$
this gives $\nu'\le(\ell+2+1/\beta)r-\ell+\log r$; with $K>\mu-2\beta$ and $\theta>(1-\delta)/\beta$ it gives
$\nu'\ge-2-\delta(\ell-2)+\log(1-\delta)$ and $\nu\ge(\ell-2)(1-\delta)$.

\emph{Certification.} Each of the resulting expressions, and $G$, is a function of $x=1/\beta$, $\ell x$, $\ell^2x$ and $1/\ell$. An
adaptive ball-arithmetic bisection over $\ell\in[\ell_A,60]$ ($5486$ boxes), together with one tail box for $\ell\ge60$ (where
$x,\ell x,\ell^2x\le3600e^{-60}$ since $\ell^ae^{-\ell}$ is decreasing for $\ell>a$), certifies $G\ge1.0092$ and the stated
constants for every $\beta\ge\beta_A$ (Table~\ref{tab:certs}, C1). No monotonicity in $\beta$ is used.
\end{proof}

The same certification gives the constants used for the other regions:
\begin{equation}\label{eq:CWconst}
  \nu'\le\nu_{CL}:=-1.95361\ \text{ for }k\le\mu-2\beta,\qquad \nu'\ge\nu_{CR}:=1.04305\ \text{ for }k\ge\mu+2\beta,
\end{equation}
and the edge values $\nu'\le-1.906674$ at the left end of $W$ and $\nu'\ge1.04305$ at the right end. (For~\eqref{eq:CWconst} one uses
that $\theta_k$ is increasing in $k$, since $S_k(t)$ increases with $k$ while $m=n-k$ decreases.)

\subsection{Riemann sums}
\begin{lemma}\label{lem:riemann}
Let $f:[0,\infty)\to\RR$ have $\TV(f)<\infty$. Then for $\theta>0$ and $k\ge1$,
$\bigl|\theta\sum_{i\le k}f(i\theta)-\int_0^{k\theta}f\bigr|\le\theta\Var_{[0,k\theta]}f\le\theta\TV(f)$.
\end{lemma}

\begin{proof}
$\theta f(i\theta)-\int_{(i-1)\theta}^{i\theta}f=\int_{(i-1)\theta}^{i\theta}(f(i\theta)-f(u))\dd u$, and $|f(i\theta)-f(u)|\le\Var_{[(i-1)\theta,i\theta]}f$.
Summing over $i$ and using additivity of the variation gives the claim.
\end{proof}

\begin{lemma}\label{lem:TV}
The total variations of $f_1,\dots,f_6$ on $[0,\infty)$ satisfy
\[
  \TV(f_1),\dots,\TV(f_6)\ \le\ 1.000001,\ \ 1.000001,\ \ 2.000001,\ \ 6.086765,\ \ 25.691943,\ \ 138.425 .
\]
\end{lemma}

\begin{proof}
On $[60,\infty)$ each $f_r$ is decreasing to $0$ (each term $l^{r-1}u^re^{-lu}$ decreases for $u>r$), so the variation there is $f_r(60)$.
On $[0,60]$ we use $6000$ cells. On a cell where ball arithmetic certifies that $f_r'$ has constant sign, $f_r$ is monotone and the
variation is $|f_r(b)-f_r(a)|$. Otherwise the cell is bisected (to depth $14$), and on the final cells the variation is bounded by
$(b-a)\sup|f_r'|\ge\int_a^b|f_r'|$. For $u\le2.5$, $f_r$ and $f_r'$ are evaluated from their Bernoulli-number expansions with a rigorous
tail bound ($|B_j|/j!\le4/(2\pi)^j$); for $u>2.5$, from the polylogarithm and the identity $uf_r'=rf_r-f_{r+1}$
(Table~\ref{tab:certs}, C2).
\end{proof}

\begin{corollary}\label{cor:sr}
$s_r\in F_r(\nu-\theta)\pm\theta\,\TV(f_r)$ and $0\le s_r\le r!\zeta_2+\theta\TV(f_r)$. In particular
$b\ge F_2(\nu-\theta)-\theta\TV(f_2)$.
\end{corollary}

\begin{proof}
By~\eqref{eq:scaled}, $s_r=\theta\sum_{i\le k}f_r(i\theta)$; apply Lemmas~\ref{lem:riemann} and~\ref{lem:Fr} with $k\theta=\nu-\theta$.
\end{proof}

Note that $f_r(\nu)=c_r\theta^r$ and $c_1=1+\nu e^{-\nu'}/(1-q)$ are \emph{exact} functions of $(\theta,\nu')$. Thus every quantity in
Proposition~\ref{prop:eps} is enclosed, given $(\theta,\nu')$, by the enclosures of Corollary~\ref{cor:sr}, and $M_\chi$ by
Corollary~\ref{cor:cstar}.

\subsection{Proof of Theorem A}\label{sec:proofA}
By Lemma~\ref{lem:W}, for $k\in W$ we have $\theta\le\theta_A$, $\nu'\in[-2.52,2.10]$ and $\nu\ge3.1158$; in particular
Corollary~\ref{cor:cstar} applies. Since $\theta q=\theta^2e^{-\nu'}$, it suffices to bound $\Gamma:=\varepsilon_k/(\theta^2e^{-\nu'})$.

\emph{Region I: $0.002\le\theta\le\theta_A$.} We cover $\theta$ by geometric cells of ratio $1.01$ and $\nu'$ by cells of width $0.01$
(clipped to $\nu\ge3.1158$). On each box we evaluate $\varepsilon$ of Proposition~\ref{prop:eps} in ball arithmetic, with $s_r$ ranging over
the enclosure of Corollary~\ref{cor:sr} for every $(\theta,\nu')$ in the box, $f_r(\nu)$ and $c_1$ exact, and
$M_\chi\le e^{-0.15926/\theta}$ (a rounded-down value, valid since $0.15926\le c_\ast$); the code asserts $\delta_0<1$, $N_0-E_D>0$ and $\Delta<M_2$ on every box, and divides by
$\min\theta^2e^{-\nu'}$ over the box. The maximum is $\Gamma\le0.42024$, attained near $\nu'=2.1$ (Table~\ref{tab:certs}, A1).

\emph{Region II: $\theta\le0.002$.} Here no box cover is possible, and we use a monotone majorant in $t=\sqrt\theta$. After dividing by
$\theta q=t^2q$, every term of $\varepsilon$ is bounded by a nonnegative combination of monomials $t^P\nu^je^{-\beta_q\nu'}$ with
coefficients that are bounded uniformly through $s_r\le r!\zeta_2+0.002\TV(f_r)$, $b\ge F_2(\nu_{\min}-0.002)-0.002\TV(f_2)$ and
$e^{-\nu}\le e^{-\nu_{\min}}$. At fixed $\nu'$ we have $\nu=\nu'+2\log(1/t)$ and
\[
  \frac{d}{dt}\bigl[t^P\nu^j\bigr]=t^{P-1}\nu^{j-1}(P\nu-2j),
\]
which is nonnegative when $P>0$ and $\nu\ge2j/P$ (and trivially when $j<0$). The program asserts, for every
monomial, either $P=j=0$ (a constant term) or $P>0$ and $\nu_{\min}>2j/P$, and evaluates the majorant at $\theta=0.002$. The Gaussian-tail and minor-arc factors $e^{-0.04b/\theta}\theta^{-2}$ and
$e^{-c_\ast/\theta}\theta^{-3.5}$ are increasing on the relevant range because $\theta<0.02b$ and $\theta<c_\ast/3.5$, which is asserted. The
result is $\Gamma\le0.17605$ for all $\theta\le0.002$ and $\nu'\in[-2.52,2.10]$ (Table~\ref{tab:certs}, A2). The details of the majorant
are in Appendix~\ref{app:regionII}. \qed

%======================================================================
\section{The step of the main term}\label{sec:B}
%======================================================================

In this section $k$ and $k+1$ both lie in $W$, and primes refer to the index $k+1$: $\theta'=\theta_{k+1}$ is the saddle point for
$(m-1,k+1)$, $q'=e^{-(K+1)\theta'}$, $\nu_{k+1}=(K+1)\theta'$, and so on.

\begin{theorem}[Theorem B]\label{thm:B}
Let $n\ge10^5$ and $k,k+1\in W$. Then
\[
  L_{2,k}-L_{2,k+1}\ \ge\ \Sigma_{\min}\,\theta_kq_k,\qquad \Sigma_{\min}=0.91918,
\]
and $\theta_{k+1}q_{k+1}\le1.01\,\theta_kq_k$. In particular $L_2$ is strictly decreasing on $W$ and
\[
  L_{2,k}-L_{2,k+1}\ \ge\ 0.9100\,\max(\theta_kq_k,\theta_{k+1}q_{k+1})\ >\ 2\Gamma_{\max}\max(\theta_kq_k,\theta_{k+1}q_{k+1}) .
\]
\end{theorem}

\subsection{Decomposition}
\begin{equation}\label{eq:stepdec}
  L_{2,k}-L_{2,k+1}=(\theta'-\theta)+\log\frac{1-q'}{1-q}-\bigl(\log M_2'-\log M_2\bigr).
\end{equation}
We shall see that $\theta'>\theta$, so the first term is positive and $q'=e^{-(K+1)\theta'}<qe^{-\theta}$. Hence
\[
  \log\frac{1-q'}{1-q}\ \ge\ T_2:=\log\Bigl(1+\frac{q(1-e^{-\theta})}{1-q}\Bigr),
\]
which is $\theta q(1+O(\theta+q))$, and the step is at least $T_2-|\log M_2'-\log M_2|$. The work is to bound $\log M_2'-\log M_2$, for which we
need to control how $(\theta,\nu,b,s_3,s_4,s_5)$ moves from $k$ to $k+1$.

\subsection{Localising the next saddle point}
Let $B_k(t):=-S_k'(t)=t^{-3}b(k,t)$ with $b(k,t):=t\sum_{i\le k}f_2(it)$. Note $S_{k+1}(t)-S_k(t)=K/(e^{Kt}-1)$, which equals $c_1-1$
at $t=\theta$. Fix $\tau=0.01$ and let $b_{\min}:=F_2(\nu-\theta)-(1+\tau)\theta\TV(f_2)$. By Corollary~\ref{cor:sr} and monotonicity of
$F_2$, $b(k,t)\ge b_{\min}$ for $t\in[\theta,(1+\tau)\theta]$.

\begin{lemma}\label{lem:B1}
\begin{enumerate}[label=(\roman*)]
\item If $\tau b_{\min}>(1+\tau)^3c_1\theta^2$, then $\theta<\theta'\le(1+\tau)\theta$.
\item Under (i), $d\theta:=\theta'-\theta\le(1+\tau)^3c_1\theta^3/b_{\min}$.
\item $d\nu:=\nu_{k+1}-\nu=\theta'+K\,d\theta\in\bigl(0,(1+\tau)\theta+(\nu/\theta)\,d\theta\bigr]$.
\item For $r=2,\dots,5$, with $\bar s_r:=r!\zeta_2+(1+\tau)\theta\TV(f_r)$,
\[
  |s_r(k+1,\theta')-s_r(k,\theta)|\le d\theta\,\frac{(1+r)\bar s_r+\bar s_{r+1}}{\theta}+(1+\tau)\theta\sup_{u\in[\nu,\nu+d\nu]}f_r(u).
\]
\end{enumerate}
\end{lemma}

\begin{proof}
(i) $S_{k+1}(\theta)=m+c_1-1>m-1=S_{k+1}(\theta')$, so $\theta'>\theta$. On $[\theta,(1+\tau)\theta]$,
$-S_{k+1}'\ge B_k(t)\ge b_{\min}/((1+\tau)\theta)^3$, so $S_{k+1}((1+\tau)\theta)\le m+c_1-1-\tau b_{\min}/((1+\tau)^3\theta^2)<m-1$.
(ii) By the mean value theorem $c_1=S_{k+1}(\theta)-S_{k+1}(\theta')=B_{k+1}(\xi)\,d\theta\ge b_{\min}\,d\theta/((1+\tau)\theta)^3$.
(iii) $\nu_{k+1}-\nu=(K+1)\theta'-K\theta$ and $K=\nu/\theta$.
(iv) $s_r(k+1,\theta')=s_r(k,\theta')+\theta'f_r(K\theta')$ with $K\theta'\in[\nu,\nu+d\nu]$. From $uf_r'=rf_r-f_{r+1}$ we get
$\partial_ts_r(k,t)=((1+r)s_r(k,t)-s_{r+1}(k,t))/t$, and $0\le s_r(k,t)\le\bar s_r$ for $t\le(1+\tau)\theta$.
\end{proof}

\subsection{Proof of Theorem B}
\emph{Region I ($0.002\le\theta\le\theta_A$).} We cover $(\theta,\nu')$ by $8526$ boxes ($\theta$-ratio $1.02$, $\nu'$-width $0.02$, clipped to
$\nu\ge3.1158$). On each box the program asserts the hypothesis of Lemma~\ref{lem:B1}(i) and encloses $d\theta$, $d\nu$ and the
increments of $s_r$. It then forms a product box $\mathcal H\subset\RR^6$ in the variables $(\theta,\nu,b,s_3,s_4,s_5)$:
$\theta\in[\theta,(1+\tau)\theta]$, $\nu\in[\nu,\nu+d\nu]$ and
$s_r\in[F_r(\nu-\theta)-(1+\tau)\theta\TV_r,\ F_r(\nu+d\nu)+(1+\tau)\theta\TV_r]$. Both the point for $k$ and the point for $k+1$ lie in
$\mathcal H$ (Corollary~\ref{cor:sr} and Lemma~\ref{lem:B1}), and $\mathcal H$ is convex. The program checks $M_2>0$ on $\mathcal H$ and
encloses $\nabla M_2$ on $\mathcal H$ by forward-mode automatic differentiation in ball arithmetic, so by the mean value theorem along the
segment
\[
  |\log M_2'-\log M_2|\le\sum_j\sup_{\mathcal H}\Bigl|\frac{\partial_jM_2}{M_2}\Bigr|\,|v_j'-v_j|,
\]
where $v=(\theta,\nu,b,s_3,s_4,s_5)$ and $v'$ is the same vector at $k+1$, with $|v_j'-v_j|$ bounded by Lemma~\ref{lem:B1}. Together with $T_2/(\theta q)$ evaluated on the box this gives
$\Sigma:=(\text{step})/(\theta q)\ge0.91918$ on every box (Table~\ref{tab:certs}, B1).

\emph{Region II ($\theta\le0.002$).} The same graded-monomial method as in Section~\ref{sec:proofA} applies. At both points every
quantity is majorised by nonnegative combinations of $t^pq^{b_q}\nu^j$, using $\theta'^e\le(1+\tau)^e\theta^e$ ($e>0$), $q'\le q$,
$\nu_{k+1}\le\rho\nu$ with an explicit $\rho$, $f_r(u)\le(\rho\nu)^rqP_{r-1}(y)/(1-y)^r$ ($y=e^{-\nu_{\min}}$, $P_{r-1}$ the Eulerian
polynomial), and $|h'(u)|\le(u+1)e^{-u}/(1-e^{-u})^2$ for $c_1=1+h(\nu)/\theta$, $h(u)=u/(e^u-1)$. Products and quotients are
handled by $|x'y'-xy|\le|x'-x|\sup|y|+\sup|x|\,|y'-y|$ and $|1/N'-1/N|\le|N'-N|/N_{\min}^2$ (primes denoting values at $k+1$), and the main term by
$T_2/(\theta q)\ge(1-\theta/2)\bigl(1-\theta q/(2(1-q))\bigr)$. The result is $\Sigma\ge0.96039$ for all $\theta\le0.002$ and
$\nu'\in[-2.52,2.10]$, and also $|\log M_2|/\theta\le0.14215$ (Table~\ref{tab:certs}, B2).

Finally $\theta'q'\le1.01\,\theta q$ because $\theta'\le1.01\theta$ and $q'<q$, and $0.91918/1.01>0.9100>0.8564=2\Gamma_{\max}$. \qed

\subsection{Consequences in the window}
Let $k^\ast$ be the largest $k\in W$ with $L_{2,k}>0$, and $k^\ast:=\min(W\cap\ZZ)-1$ if there is no such $k$ (so $L_{2,k}\le0$ for every $k\in W$ with $k>k^\ast$). Call $k\in W$
\emph{undetermined} if $|L_{2,k}|\le\varepsilon_k$. If $k$ is not undetermined then $\operatorname{sign}L_k=\operatorname{sign}L_{2,k}$
by Theorem~\ref{thm:A}.

\begin{corollary}\label{cor:onek}
Let $n\ge10^5$. At most one $k\in W$ is undetermined, and if there is one it belongs to $\{k^\ast,k^\ast+1\}$. Hence the signs of
$L_k$ for $k\in W$, read left to right, have the form $+\cdots+-\cdots-$ (either block may be empty).
\end{corollary}

\begin{proof}
For $k\le k^\ast-1$, Theorem~\ref{thm:B} gives $L_{2,k}\ge L_{2,k}-L_{2,k+1}>0.91\,\theta_kq_k>\Gamma_{\max}\theta_kq_k\ge\varepsilon_k$.
For $k\ge k^\ast+2$, $L_{2,k}\le L_{2,k}-L_{2,k-1}\le-0.91\,\theta_kq_k$. So only $k^\ast$ and $k^\ast+1$ can be undetermined. If both
were, then
\[
  L_{2,k^\ast}-L_{2,k^\ast+1}\le\varepsilon_{k^\ast}+\varepsilon_{k^\ast+1}\le2\Gamma_{\max}\max(\theta q)<0.91\max(\theta q),
\]
contradicting Theorem~\ref{thm:B}. The determined signs are $+$ for $k\le k^\ast$ and $-$ for $k\ge k^\ast+1$, and a single undetermined
index adjacent to the change cannot create a second change.
\end{proof}

\begin{remark}[Location of the mode]
Using the edge values of $\nu'$ after Lemma~\ref{lem:W} one finds $L_2/\theta\ge5.588$ at the left end of $W$ and
$L_2/\theta\le-0.505$ at the right end (Table~\ref{tab:certs}, B3), so the sign change of $L_2$, and hence the mode of the row, lies
inside $W$. This is not needed for unimodality.
\end{remark}

%======================================================================
\section{The regions D, CL, CR and E}\label{sec:CLCR}
%======================================================================

\subsection{A zeroth-order estimate}
Outside the window a much cruder version of Proposition~\ref{prop:eps} suffices. Let $c_0>0$ and $\varphi_0=c_0\theta$.

\begin{lemma}[Lemma C0]\label{lem:C0}
Suppose $\re\psi(\varphi)\le-(1-\delta)B\varphi^2/2$ for $|\varphi|\le\varphi_0$, with $0\le\delta<1$. Put
\[
  T_1=\frac{c_1\kappa_3}{2B^2},\quad T_2=\frac{c_2+c_1^2}{2B},\quad A_3=\kappa_3B^{-3/2},\quad w_0=\varphi_0\sqrt B,\quad
  \mathcal M=\sqrt{\frac B{2\pi}}\int_{\varphi_0\le|\varphi|\le\pi}|\chi| .
\]
If the denominator below is positive, then
\[
  \Bigl|\frac{R_k}{g_0}-1\Bigr|\le\rho:=\frac{T_1(1-\delta)^{-5/2}+T_2(1-\delta)^{-3/2}+2\mathcal M}
  {1-\erfc(w_0/\sqrt2)-\sqrt{2/\pi}\,A_3/\bigl(3(1-\delta)^2\bigr)-\mathcal M}.
\]
Moreover $\mathcal M\le\sqrt{2\pi B}\,\sup_{\varphi_0\le|\varphi|\le\pi}|\chi|$.
\end{lemma}

\begin{proof}
We use~\eqref{eq:rep2} and measure all integrals in units of $\sqrt{2\pi/B}$.
\emph{Numerator, central arc.} Since $(e^H)''=(H''+H'^2)e^H$ and $\re H\le0$, Taylor's theorem gives $e^H-1=\ii c_1\varphi+E(\varphi)$ with
$|E|\le(c_2+c_1^2)\varphi^2/2$. By $\chi(-\varphi)=\overline{\chi(\varphi)}$, $\int_{|\varphi|\le\varphi_0}\chi\varphi=2\ii\int_0^{\varphi_0}\varphi\im\chi$,
and $|\im\chi|=e^{\re\psi}|\sin\im\psi|\le e^{-(1-\delta)B\varphi^2/2}\kappa_3|\varphi|^3/6$ by Corollary~\ref{cor:omega}. Integrating
$\varphi^4$ and $\varphi^2$ against $e^{-(1-\delta)B\varphi^2/2}$ over $\RR$ gives the terms $T_1(1-\delta)^{-5/2}$ and $T_2(1-\delta)^{-3/2}$.
\emph{Numerator, minor arcs.} $|e^H-1|\le2$ gives $2\mathcal M$.
\emph{Denominator.} $D$ is real and positive. On $|\varphi|\le\varphi_0$ write $\chi=e^{-B\varphi^2/2}e^\omega$; by
Corollary~\ref{cor:omega} and~\ref{F1}, $|e^\omega-1|\le|\omega|e^{\delta B\varphi^2/2}\le\frac{\kappa_3|\varphi|^3}6e^{\delta B\varphi^2/2}$, whose
weighted integral is $\sqrt{2/\pi}A_3/(3(1-\delta)^2)$. Also $\int_{|\varphi|\le\varphi_0}e^{-B\varphi^2/2}=\sqrt{2\pi/B}\,(1-\erfc(w_0/\sqrt2))$,
and the minor arcs contribute at most $\mathcal M$.
\end{proof}

\begin{lemma}[Admissible $\delta$]\label{lem:delta}
Both $\delta_T:=\kappa_4\varphi_0^2/(12B)$ and, when $c_0k\theta<2$,
$\delta_C:=1-\bigl(1-(c_0k\theta)^2/12\bigr)\log(1+c_0^2)/c_0^2$ satisfy the hypothesis of Lemma~\ref{lem:C0} (when they are $<1$).
\end{lemma}

\begin{proof}
For $\delta_T$ use Corollary~\ref{cor:omega}: $\re\psi\le-B\varphi^2/2+\kappa_4\varphi_0^2\varphi^2/24$. For $\delta_C$, by~\eqref{eq:modchi},
$-2\re\psi=\sum_i\log(1+\sin^2(i\varphi/2)/\sinh^2(i\theta/2))$. For $i\le k$ and $|\varphi|\le c_0\theta$ the angle $y=i|\varphi|/2$ is
$<1$, and $\sin^2y\ge(y-y^3/6)^2\ge y^2(1-y^2/3)\ge\lambda y^2$ with $\lambda:=1-(c_0k\theta)^2/12$. With
$a_i:=(i\varphi/2)^2/\sinh^2(i\theta/2)\le c_0^2$, two applications of~\ref{F2} give
$\log(1+\lambda a_i)\ge\lambda\log(1+a_i)\ge\lambda a_i\log(1+c_0^2)/c_0^2$, and $\sum_ia_i=B\varphi^2$.
\end{proof}

From $\log R_k=-\theta-\log(1-q)+\log(R_k/g_0)$, $q\le-\log(1-q)\le q/(1-q)$ and $1-q\le\nu$ we obtain, whenever $\rho<1$,
\begin{equation}\label{eq:C0cons}
  \log R_k\le-\theta+\frac q{1-q}+\rho,\qquad \log R_k\ge-\theta+q+\log(1-\rho),\qquad \log R_k\ge-\theta-\log\nu+\log(1-\rho).
\end{equation}

\subsection{A bound uniform in \texorpdfstring{$\theta$}{theta}}
\begin{lemma}[Lemma C1]\label{lem:C1}
Fix $0<v_1<v_2\le\infty$, $c_0>0$ and $\bar\theta>0$. Suppose $\Psi>5\bar\theta$ is such that $\theta S(\varphi)\ge\Psi$ for
$c_0\theta\le|\varphi|\le\pi$ whenever $\theta\le\bar\theta$ and $k\theta\in[v_1,v_2]$. Then there is an explicit number
$X=X(\bar\theta,v_1,v_2,c_0)$ such that $\rho\le\theta X$ for every saddle point $\theta\le\bar\theta$ with $k\theta\in[v_1,v_2]$.
\end{lemma}

\begin{proof}
Write $v=k\theta$. We bound the ingredients of Lemma~\ref{lem:C0}:
\begin{itemize}
\item $s_r\le\min\bigl(F_r(v_2)+\theta\TV(f_r),\ v\sup f_r\bigr)$ for $r=3,4$ (Corollary~\ref{cor:sr}, and each term $\theta f_r(i\theta)\le\theta\sup f_r$);
  with $\sup f_3\le2.0302009$, $\sup f_4\le6.1253939$;
\item $b\ge\max\bigl(F_2(v_1)-\theta,\ v_1f_2(v_2)\bigr)$ and $b\le F_2(v_2)$, by comparing the decreasing $f_2\le1$ with its integral;
\item $c_1\theta=\theta+f_1(\nu)\le\theta+f_1(v_1)$ and $c_2\theta^2=f_2(\nu)\le f_2(v_1)$, as $\nu\ge v\ge v_1$;
\item $\sup|\chi|\le e^{-\Psi/(2\theta)}$, hence $\mathcal M\le\sqrt{2\pi b}\,\theta^{-3/2}e^{-\Psi/(2\theta)}$.
\end{itemize}
In scaled form $T_1=\theta(c_1\theta)s_3/(2b^2)$, $T_2=\theta(c_2\theta^2+(c_1\theta)^2)/(2b)$, $A_3=s_3\theta^{1/2}b^{-3/2}$,
$w_0=c_0(b/\theta)^{1/2}$ and $\delta_T=c_0^2s_4/(12b)$. The quantity $\rho$ is increasing in $T_1,T_2,\mathcal M,\delta,A_3$ and
$\erfc(w_0/\sqrt2)$ (its numerator increases and its denominator decreases in each). After substituting the bounds above,
$T_1/\theta$, $T_2/\theta$, $A_3$ and $\delta$ are nondecreasing in $\theta$ and $w_0$ is nonincreasing; this is checked by computing
each derivative symbolically as a polynomial with nonnegative coefficients over a positive denominator (Table~\ref{tab:certs}, M).
Finally $\frac{d}{d\theta}\log(\mathcal M/\theta)=(\Psi-5\theta)/(2\theta^2)>0$ for $\theta<\Psi/5$. Hence $\rho/\theta$ is bounded by its
value with all ingredients taken at $\bar\theta$, which defines $X$. (The program takes the best of $c_0\in\{0.4,0.6,0.8,1,1.25,1.5,2\}$
and of $\delta_T,\delta_C$.)
\end{proof}

In this section $\Psi$ is provided by Lemma~\ref{lem:minor} with $u_{\mathrm{top}}=v_1$ (since $(k+1)\theta\ge k\theta\ge v_1$).

\subsection{The right tail CR}
\begin{proposition}\label{prop:CR}
Let $n\ge10^5$ and $\mu+2\beta\le k\le\lceil n/2\rceil$. Then $\log R_k\le-0.614\,\theta<0$.
\end{proposition}

\begin{proof}
Since $m\ge n-\lceil n/2\rceil\ge(n-1)/2$ and $\theta^2m\le\zeta_2$, $\theta\le\theta_{CR}:=\sqrt{2\zeta_2/(n-1)}\le0.0057358$. By~\eqref{eq:CWconst},
$\nu'\ge\nu_{CR}=1.04305$, so $v=\nu'-\log\theta-\theta\ge1.04305+\log(1/\theta_{CR})-\theta_{CR}\ge6.198$; both bounds improve with $n$.
Lemma~\ref{lem:minor} with $\theta_{\max}=\theta_{CR}$ and $u_{\mathrm{top}}=6.198$ gives $\Psi=0.33808>5\theta_{CR}$ (\script{ws\_mech\_cd}), and Lemma~\ref{lem:C1} on the single cell
$[6.198,\infty)$ (where only $F_r(\infty)=r!\zeta_2$ is used) gives $X\le0.0321$. By~\eqref{eq:C0cons}, with $q/\theta=e^{-\nu'}$ and $q\le e^{-v}$,
\[
  \log R_k\le\theta\Bigl(\frac{e^{-\nu'}}{1-q}+X-1\Bigr)\le\theta(0.3852-1).
\]
(Table~\ref{tab:certs}, R1.)
\end{proof}

\subsection{The left tail CL}
\begin{proposition}[CL-a]\label{prop:CLa}
For every $n$, if $47\le k\le n-2$ and $k\theta\le0.1$, then $R_k>1$.
\end{proposition}

\begin{proof}
Put $K_0=47$, $V_0=0.1$, $\theta_0=V_0/K_0$, so $\theta\le\theta_0$. Here $b\ge vf_2(V_0)$, $s_r\le v\sup f_r$, $c_1\theta\le1+\theta_0$ and
$f_2\le1$, which give $T_1\le(1+\theta_0)\sup f_3/(2K_0f_2(V_0)^2)$, $T_2\le(1+(1+\theta_0)^2)/(2K_0f_2(V_0))$,
$A_3\le\sup f_3\,K_0^{-1/2}f_2(V_0)^{-3/2}$ and $w_0\ge c_0(K_0f_2(V_0))^{1/2}$. We bound $\mathcal M\le\mathcal N_1+\mathcal N_2+\mathcal N_{34}$
according to $\varphi\in[c_0\theta,10\theta]$, $[10\theta,\pi/k]$ and $[\pi/k,\pi]$; since $10\theta=10v/k<\pi/k$ these cover the minor arcs.

\emph{$\mathcal N_1$.} Split $c=\varphi/\theta\in[c_0,10]$ into $60$ cells $[c_j,c_{j+1}]$. As $i\varphi/2\le c_{j+1}V_0/2=:T_j<\pi/2$,
$\sin(i\varphi/2)\ge\sigma_j\,i\varphi/2$ with $\sigma_j=\sin T_j/T_j$, and by~\ref{F2} each summand of $S$ is at least
$\lambda_jc^2f_2(i\theta)$ with $\lambda_j=\log(1+\sigma_j^2c_{j+1}^2)/c_{j+1}^2$. Since $\sum_ic^2f_2(i\theta)=B\varphi^2$, this gives
$|\chi|\le e^{-\lambda_jB\varphi^2/2}$, and
$2\sqrt{B/2\pi}\int_{c_j\theta}^\infty e^{-\lambda_jB\varphi^2/2}\dd\varphi=\erfc\bigl(c_j\theta\sqrt{\lambda_jB/2}\bigr)/\sqrt{\lambda_j}$ with
$\theta^2B=b/\theta\ge K_0f_2(V_0)$.

\emph{$\mathcal N_2$.} Here every angle $i\varphi/2$ is at most $\pi/2$, hence $\sin(i\varphi/2)\ge i\varphi/\pi$, and with
$c=\varphi/\theta$ each summand of $S$ is at least $\log\bigl(4c^2f_2(V_0)/\pi^2\bigr)$. Hence, using $\theta\sqrt B\le\sqrt k$,
\[
  |\chi|\le\Bigl(\frac{4f_2(V_0)}{\pi^2}\Bigr)^{-k/2}c^{-k},\qquad
  \mathcal N_2\le\sqrt{\frac2\pi}\,\sqrt k\,\Bigl(\frac{4f_2(V_0)}{\pi^2}\Bigr)^{-k/2}\frac{10^{1-k}}{k-1}.
\]

\emph{$\mathcal N_{34}$.} Put $d=\varphi/2\in[\pi/(2k),\pi/2]$. By~\ref{F2} each summand is at least $\sin^2(id)w_i$ with
$w_i=\log(1+1/\sinh^2(i\theta/2))\ge w_k$, and by~\eqref{eq:dirichlet} $\sum_{i\le k}\sin^2(id)\ge k/2-1/(2\sin d)$. Since
$\sin d\ge(\pi/2k)\sinc(\pi/(2K_0))$, this is at least $2gk$ with $g:=\frac12\bigl(\frac12-1/(\pi\sinc(\pi/2K_0))\bigr)>0$. Also
$w_k\ge-2\log\sinh(v/2)\ge-2\log(v\sigma_0/2)$ with $\sigma_0=\sinh(V_0/2)/(V_0/2)$. So $|\chi|\le(v\sigma_0/2)^{2gk}$, and with
$B\le k^3/v^2$, $\mathcal N_{34}\le\sqrt{2\pi}\,k^{3/2}v^{-1}(v\sigma_0/2)^{2gk}$, which is increasing in $v$ and decreasing in $k$ (both
asserted), so it is evaluated at $(K_0,V_0)$.

With $c_0=0.6$ and $\delta=\min(\delta_T,\delta_C)$ this gives $\mathcal N_1\le3.7\cdot10^{-4}$, $\mathcal N_2\le10^{-37}$,
$\mathcal N_{34}\le6.4\cdot10^{-8}$ and $\rho\le0.0673$. By the third inequality of~\eqref{eq:C0cons}, $\log R_k>0$ as soon as
$1-\rho>e^\theta\nu$, and $e^\theta\nu\le e^{\theta_0}V_0(1+1/K_0)\le0.1024$ (Table~\ref{tab:certs}, R2).
\end{proof}

\begin{proposition}[CL-b]\label{prop:CLb}
Let $n\ge10^5$ and $k_D<k\le\mu-2\beta$ with $k\theta\ge0.1$. Then $R_k>1$.
\end{proposition}

\begin{proof}
Here $k\ge79$, $\theta\le\theta_A$ and $\nu'\le\nu_{CL}=-1.95361$ by~\eqref{eq:CWconst}. Cover $v=k\theta\in[0.1,\infty)$ by the $32$ geometric
cells $[v_1,v_2]$ of ratio $1.15$ starting at $0.1$ and the final cell $[0.1\cdot1.15^{32},\infty)=[8.7565\ldots,\infty)$. On each cell $\theta\le\bar\theta:=\min(\theta_A,v_2/78)$; let
$X$ be the constant of Lemma~\ref{lem:C1}. By~\eqref{eq:C0cons}, $\log R_k\ge-\theta+q-\theta X/(1-\theta X)$, which is positive as soon as
$q/\theta>1+X/(1-\bar\theta X)$. For $q/\theta$ we use the larger of $e^{-\nu'}\ge e^{1.95361}$ and $e^{-(v_2+\bar\theta)}/\bar\theta$. Every
cell passes, with minimum ratio $5.82$ (Table~\ref{tab:certs}, R3).
\end{proof}

\subsection{Small \texorpdfstring{$k$}{k}: the region D}\label{sec:D}
\begin{lemma}\label{lem:D1}
For $1\le k\le n$, $\binom{n-1}{k-1}\le k!\,p(n,k)\le\binom{n-1+k(k-1)/2}{k-1}$.
\end{lemma}

\begin{proof}
Sorting a composition of $n$ into $k$ positive parts gives a partition into $k$ parts, and each partition arises from at most $k!$
compositions. Conversely $\lambda\mapsto(\lambda_i+k-i)_{i\le k}$ maps partitions of $n$ into $k$ parts injectively to partitions of
$n+k(k-1)/2$ into $k$ distinct parts, each of which arises from exactly $k!$ compositions.
\end{proof}

\begin{proposition}\label{prop:D}
Let $n\ge10^5$ and $1\le k\le k_D(n):=\lfloor1.7n^{1/3}\rfloor$. Then $p(n,k+1)>p(n,k)$.
\end{proposition}

\begin{proof}
Let $\Phi(k,n):=\log(n-k)-\log(k(k+1))-k(k-1)^2/(2(n-k+1))$. \emph{Step 1: $\Phi>0$ implies the claim.} By Lemma~\ref{lem:D1} it suffices
that $\binom{n-1}{k}/\binom{N-1}{k-1}>k+1$ with $N=n+D$, $D=k(k-1)/2$. The left side equals
$\frac{n-k}k\prod_{j=0}^{k-2}(1+D/(n-1-j))^{-1}\ge\frac{n-k}k\exp(-(k-1)D/(n-k+1))$, by $1+y\le e^y$ and $n-1-j\ge n-k+1$.

\emph{Step 2: monotonicity in $k$.} For real $1\le k<n$,
\[
  \partial_k\Phi=-\frac1{n-k}-\frac1k-\frac1{k+1}-\frac{(k-1)\bigl[(3k-1)(n-k+1)+k(k-1)\bigr]}{2(n-k+1)^2}<0 .
\]
Put $\gamma=1.7$ and $s=n^{1/3}\ge s_A:=10^{5/3}$; since $\gamma s<n-1$, $\Phi(k,n)\ge\Phi(\gamma s,s^3)$ for $k\le\gamma s$.

\emph{Step 3: reduction to one variable.} Using $k(k-1)^2/(n-k+1)\le k^3/(n-k)$, which follows from
$k^3(n-k+1)-k(k-1)^2(n-k)=k[(n-k)(2k-1)+k^2]\ge0$, we get $\Phi(\gamma s,s^3)\ge\Psi_D(s)$, where
\[
  \Psi_D(s):=\log\frac{s^2-\gamma}{\gamma(\gamma s+1)}-\frac{\gamma^3}{2(1-\gamma/s^2)},
\]
and, exactly,
\[
  \Psi_D'(s)=\frac{\gamma s^2+2s+\gamma^2}{(s^2-\gamma)(\gamma s+1)}+\frac{\gamma^4s}{(s^2-\gamma)^2}>0\quad(s>\sqrt\gamma).
\]
\emph{Step 4.} $\Psi_D(s_A)=0.30456\ldots>0$ (ball arithmetic), so $\Psi_D(s)>0$ for all $s\ge s_A$ (Table~\ref{tab:certs}, R4).
\end{proof}

\subsection{Large \texorpdfstring{$k$}{k}: the region E}
\begin{lemma}\label{lem:E}
If $2k\ge n$ then $p(n,k)=p(n-k)$. Consequently $p(n,k+1)\le p(n,k)$ for $\lceil n/2\rceil\le k\le n-1$, with equality only for $k=n-1$.
\end{lemma}

\begin{proof}
Subtracting $1$ from each part is a bijection from partitions of $n$ into exactly $k$ parts to partitions of $n-k$ into at most $k$ parts,
and the latter condition is vacuous when $n-k\le k$. For $k\ge\lceil n/2\rceil$ also $2(k+1)\ge n$, so $p(n,k+1)=p(j)$ and $p(n,k)=p(j+1)$
with $j=n-k-1$. Appending a part $1$ injects partitions of $j$ into those of $j+1$, and misses $(j+1)$ when $j\ge1$.
\end{proof}

%======================================================================
\section{Proof of the Main Theorem for \texorpdfstring{$n\ge10^5$}{n >= 10\^5}}\label{sec:assembly}
%======================================================================

Let $n\ge10^5$. The five ranges
\[
  D=[1,k_D],\ CL=(k_D,\mu-2\beta],\ W=(\mu-2\beta,\mu+2\beta),\ CR=[\mu+2\beta,\lceil n/2\rceil],\ E=[\lceil n/2\rceil,n-1]
\]
cover $[1,n-1]$, because $\mu+2\beta=\beta(\ell+2)\le n/2-1$: indeed $\beta(\ell+2)/(n/2)=(2/\zeta_2)(\ell+2)e^{-\ell}$ is decreasing for
$\ell>-1$ and is below $0.04$ at $n=10^5$. (If $k_D\ge\mu-2\beta$, the range $CL$ is empty and nothing changes.) Every $k$ in $CL\cup W\cup CR$
satisfies $79\le k\le\lceil n/2\rceil\le n-2$, so Proposition~\ref{prop:rep} applies there. The hypotheses of every lemma used hold for every
$n\ge10^5$; this bookkeeping is collected in Table~\ref{tab:thresholds} and checked by a separate program (Table~\ref{tab:certs}, T).

Now:
\begin{itemize}
\item $R_k>1$ for $k\in D$ (Proposition~\ref{prop:D});
\item $R_k>1$ for $k\in CL$: such $k$ satisfy $k\ge79\ge47$, so Proposition~\ref{prop:CLa} applies if $k\theta\le0.1$ and
  Proposition~\ref{prop:CLb} if $k\theta\ge0.1$;
\item on $W$ the signs of $\log R_k$ have the form $+\cdots+-\cdots-$ (Corollary~\ref{cor:onek});
\item $R_k<1$ for $k\in CR$ (Proposition~\ref{prop:CR}) and $R_k\le1$ for $k\in E$ (Lemma~\ref{lem:E}).
\end{itemize}
Concatenating, the signs of $p(n,k+1)-p(n,k)$, $1\le k\le n-1$, change at most once, from nonnegative to nonpositive. Hence the row is
unimodal. \qed

\begin{table}[ht]
\caption{Hypotheses of the lemmas and why they hold for all $n\ge10^5$.}\label{tab:thresholds}
\small
\begin{tabular}{@{}>{\raggedright\arraybackslash}p{0.27\textwidth}>{\raggedright\arraybackslash}p{0.33\textwidth}>{\raggedright\arraybackslash}p{0.33\textwidth}@{}}
\toprule
Lemma & Hypothesis & Holds because \\
\midrule
Cor.~\ref{cor:cstar} (minor arcs, window) & $\theta\le0.0042$, $(k+1)\theta\ge3$ & Lemma~\ref{lem:W}: $\theta\le0.0040939$, $\nu\ge3.1158$ \\
Thm.~\ref{thm:A}, Region I & $\theta\in[0.002,\theta_A]$, $\nu'\in[-2.52,2.10]$, $\nu\ge3.1158$ & Lemma~\ref{lem:W} \\
Thm.~\ref{thm:A}, Region II & $\theta\le0.002$, $\nu'\in[-2.52,2.10]$, $\theta<c_\ast/3.5$ & Lemma~\ref{lem:W}; $0.15926/3.5>0.002$ \\
Thm.~\ref{thm:B} & same ranges at $k$; hypothesis of Lemma~\ref{lem:B1}(i) & asserted on every box and cell \\
Cor.~\ref{cor:onek} & $\Sigma_{\min}/1.01>2\Gamma_{\max}$ & $0.91008>0.8564$ \\
Prop.~\ref{prop:CR} & $\theta\le\theta_{CR}$, $v\ge6.198$, $\Psi>5\theta_{CR}$ & $\theta_{CR}$ decreasing in $n$; $0.33808>0.0287$ \\
Prop.~\ref{prop:CLa} & $k\ge47$, $k\theta\le0.1$ & independent of $n$ \\
Prop.~\ref{prop:CLb} & $k\ge79$, $\theta\le\theta_A$, $\nu'\le\nu_{CL}$ & $k_D(n)\ge78$; \eqref{eq:CWconst} \\
Prop.~\ref{prop:D} & $n\ge10^5$, $k\le1.7n^{1/3}$ & $\Psi_D$ increasing \\
Prop.~\ref{prop:rep} & $2\le k\le n-2$ & $79\le k\le\lceil n/2\rceil$ on $CL\cup W\cup CR$ \\
\bottomrule
\end{tabular}
\end{table}

%======================================================================
\section{Finite verification for \texorpdfstring{$n\le2\cdot10^5$}{n <= 2*10\^5}}\label{sec:finite}
%======================================================================

By Lemma~\ref{lem:E}, the row $n$ is unimodal if and only if the signs of $p(n,k)-p(n,k-1)$ for $2\le k\le\lceil n/2\rceil+1$ contain no
$+$ after a $-$. All $p(n,k)$ with $n\le N$ are produced simultaneously by the column recurrence
\[
  v[j]\leftarrow v[j]+v[j-k]\qquad(j=k,k+1,\dots),\qquad k=1,2,\dots,
\]
started from $v=(1,0,0,\dots)$: after step $k$, $v[j]$ is the number of partitions of $j$ into parts $\le k$, which equals $p(j+k,k)$
by Lemma~\ref{lem:comb}. The cost is $O(N^2)$ additions.

Two independent programs implement this.
\begin{enumerate}
\item \emph{Exact integers} (\script{ws\_fc\_exact.c}). Fixed-width multi-limb unsigned integers ($27$ limbs of $64$ bits for $N=2\cdot10^5$,
  justified by $p(N)<e^{\pi\sqrt{2N/3}}$; a carry beyond the last limb aborts the run). Only additions and comparisons are performed.
  Since $v[j+1]$ is not yet updated when $v[j]$ has just been updated at step $k$, the comparison of $p(n,k)$ with $p(n,k-1)$ needs no
  second array. The program uses no floating point and no threads.
\item \emph{Interval arithmetic} (\script{ws\_fc\_check.c}). Two copies of the array in x87 extended precision ($64$-bit significand,
  no overflow since $\log_2p(2\cdot10^5)\approx1660$), one updated with rounding toward $-\infty$ and one toward $+\infty$. Since every
  operation is an addition of nonnegative numbers, induction gives $\mathrm{lo}[j]\le v[j]\le\mathrm{hi}[j]$ throughout. A pair is classified
  as $+$ or $-$ only if the intervals are disjoint; no pair remained unclassified.
\end{enumerate}
For $N=2\cdot10^5$ both programs report no violation; there are $256\,705\,214$ increasing pairs, $9\,743\,394\,777$ decreasing pairs and $8$
ties (all for small $n$). They write a per-$n$ log (mode and verdict) and the two logs are byte-identical. For $N=2000$ both logs also agree
byte-for-byte with a third, Python big-integer computation via $p(n,k)=p(n-1,k-1)+p(n-k,k)$. As an illustration, the mode of the row is
$k=2042$ for $n=2\cdot10^5$.

Since $2\cdot10^5>10^5$, Sections~\ref{sec:assembly} and~\ref{sec:finite} together prove the Main Theorem for every $n\ge1$.

%======================================================================
\section{The computer-assisted parts}\label{sec:computer}
%======================================================================

\subsection{Method}
All inequalities in the analytic part that involve explicit numbers are certified in Arb~\cite{Arb} ball arithmetic (python-flint 0.9.0).
A ball enclosure of a function on a box contains every value the function takes on that box, so a strict inequality certified on each box of
a cover holds on the union. Where a statement is about all $\theta$ near $0$ or all $\beta\ge\beta_A$, it is reduced to finitely many boxes
either by a tail box in variables that are monotone on the tail (Lemma~\ref{lem:W}), or by a monotone majorant whose monotonicity is checked
symbolically (Section~\ref{sec:proofA}, Lemma~\ref{lem:C1}). Floating-point grids are used only to choose the cells. Integrals are bounded
below by right Riemann sums of decreasing integrands. Exact algebraic identities (Gaussian moments, the polynomial identity of
Lemma~\ref{lem:gap}, derivatives, the Dirichlet identity~\eqref{eq:dirichlet}, the telescoping in Lemma~\ref{lem:Fr}) are checked with
SymPy.

\subsection{List of certified statements}
Table~\ref{tab:certs} lists every computer-assisted statement used in the proof, the program that certifies it (in the directory
\texttt{analytic/} or \texttt{finite/} of the repository), and the certified value. The driver script \texttt{run\_all.sh} runs all
the analytic programs and the exact finite check (to $2\cdot10^5$ with the option \texttt{--full}) and compares each output with a
stored reference; the interval program F2 is run separately (see the repository README); on a two-core machine the analytic part takes about five minutes and the
exact finite check to $2\cdot10^5$ about eight minutes.

{\small
\begin{longtable}{@{}>{\raggedright\arraybackslash}p{0.07\textwidth}>{\raggedright\arraybackslash}p{0.25\textwidth}>{\raggedright\arraybackslash}p{0.36\textwidth}>{\raggedright\arraybackslash}p{0.22\textwidth}@{}}
\caption{Computer-assisted statements. ``Sanity'' rows test statements that are proved in the text and are not part of the proof.}\label{tab:certs}\\
\toprule
Item & Program & Statement certified & Value \\
\midrule
\endfirsthead
\toprule
Item & Program & Statement certified & Value \\
\midrule
\endhead
\bottomrule
\endfoot
S1 & \script{ws\_mech\_elem} & elementary monotonicity facts; Dirichlet identity~\eqref{eq:dirichlet}; Abel summation; $uf_r'=rf_r-f_{r+1}$ & exact (SymPy) \\
S2 & \script{ws\_mech\_lemA} & Gaussian moments; $D_0=\sqrt{2\pi/B}N_0$; $\tilde N_0$ real; Lemma~\ref{lem:gap}; scaling invariance; the decompositions in Lemma~\ref{lem:central}; algebraic reduction to $G$ in Lemma~\ref{lem:W}; independent re-implementation of $\varepsilon$ agrees with the production code & exact; agreement $1.7\cdot10^{-13}$ \\
S3 & \script{ws\_close\_elem} & telescoping in Lemma~\ref{lem:Fr} and its limits; Arb quadrature of $F_r$; identities in Lemma~\ref{lem:B1}; Gaussian integrals in Lemmas~\ref{lem:C0},~\ref{lem:delta}, Prop.~\ref{prop:CLa} & exact / Arb \\
C1 & \script{ws\_mech\_window} & Lemma~\ref{lem:W} and~\eqref{eq:CWconst} for all $\beta\ge\beta_A$ & $G\ge1.0092$, $\theta_A\le0.0040939$, $\nu'\in[-2.5102,2.0840]$, $\nu\ge3.1158$, $\nu_{CL}\le-1.95361$, $\nu_{CR}\ge1.04305$ \\
C2 & \script{ws\_mech\_tv} & Lemma~\ref{lem:TV} & see Lemma~\ref{lem:TV} \\
C3 & \script{ws\_mech\_minor} & Corollary~\ref{cor:cstar} & $c_\ast\ge0.15973$ \\
A1 & \script{ws\_p12\_boxes} & Theorem~\ref{thm:A}, $\theta\in[0.002,\theta_A]$ & $\Gamma\le0.42024$ \\
A2 & \script{ws\_p12\_asym} & Theorem~\ref{thm:A}, $\theta\le0.002$ & $\Gamma\le0.17605$ \\
B1 & \script{ws\_b\_step} & Theorem~\ref{thm:B}, $\theta\in[0.002,\theta_A]$ & $\Sigma\ge0.91918$ \\
B2 & \script{ws\_b\_asym} & Theorem~\ref{thm:B}, $\theta\le0.002$ & $\Sigma\ge0.96039$ \\
B3 & \script{ws\_b\_edges}, \script{ws\_mech\_edges} & edge values of $L_2/\theta$ (not needed for unimodality) & $\ge5.588$; $\le-0.505$ \\
M & \script{ws\_mech\_mono} & monotonicity in $\theta$ of the majorants in Lemma~\ref{lem:C1} and Section~\ref{sec:proofA} & exact (SymPy) \\
R1 & \script{ws\_mech\_cd}, \script{ws\_cd\_CR} & Proposition~\ref{prop:CR} & $0.3852<1$ \\
R2 & \script{ws\_mech\_cla}, \script{ws\_cd\_CLa} & Proposition~\ref{prop:CLa} & $\rho\le0.0673<0.8976$ \\
R3 & \script{ws\_mech\_cd}, \script{ws\_cd\_CLb} & Proposition~\ref{prop:CLb} & min.\ ratio $5.82>1$ \\
R4 & \script{ws\_close\_D}, \script{ws\_cd\_D} & Proposition~\ref{prop:D}: $\partial_k\Phi$, $\Psi_D'$ in closed form, $\Psi_D(s_A)>0$ & $\Psi_D(s_A)=0.30456$ \\
T & \script{ws\_close\_thresh} & Table~\ref{tab:thresholds} and the cover in Section~\ref{sec:assembly}, reading the values certified above & --- \\
F1 & \script{ws\_fc\_exact.c} & Section~\ref{sec:finite}, exact integers, $n\le2\cdot10^5$ & 0 violations \\
F2 & \script{ws\_fc\_check.c} & Section~\ref{sec:finite}, directed-rounding intervals, $n\le2\cdot10^5$ & 0 violations, identical log \\
\midrule
Sanity & \script{ws\_close\_rep} & \eqref{eq:rep1} in exact rational arithmetic ($4959$ cases), \eqref{eq:rep2} by quadrature, Lemmas~\ref{lem:psider}, \ref{lem:Hder}, \ref{lem:D1}, \ref{lem:E} on instances & --- \\
Sanity & \script{ws\_b\_check}, \script{ws\_p12\_window}, \script{ws\_cd\_window}, \script{ws\_p12\_minor}, \script{ws\_b\_minorM} & pointwise cross-checks of the above with independent code & --- \\
\end{longtable}
}

\subsection{Trust base}
The proof relies on: the correctness of Arb's ball arithmetic as exposed by python-flint 0.9.0 (including its polylogarithm, zeta, gamma,
error-function and integration routines); SymPy 1.14 for exact symbolic identities; the GCC 13.3 compiler for the integer code of the finite
check (cross-checked by the independent interval program and, for $n\le2000$, by Python big integers); and the following classical results,
which are quoted rather than proved: Taylor's theorem with integral remainder, the mean value theorem, the Weierstrass M-test for termwise
differentiation, and additivity of total variation. All other steps are proved in this paper.

%======================================================================
\appendix
\section{The error bound in scaled form}\label{app:eps}
%======================================================================

We record the expression for $\varepsilon_k$ of Proposition~\ref{prop:eps} as it is evaluated. Set $B=1$ and replace $\kappa_r,c_r$ by
$A_r,C_r$ of~\eqref{eq:AC} and $\varphi_0$ by $w_0=0.4\sqrt{b/\theta}$. Write $a_r=A_r/r!$, $\delta_0=A_4w_0^2/12+A_6w_0^4/360$ and
$B'=1-\delta_0$. For a polynomial $P(t)=\sum p_jt^j$ with $p_j\ge0$ and $A>0$ put
\[
  \mathcal I_A[P]:=\frac1{\sqrt{2\pi}}\int_\RR P(|w|)e^{-Aw^2/2}\dd w=\sum_jp_j\,\frac{2^{j/2}\Gamma(\frac{j+1}2)}{\sqrt\pi}\,A^{-(j+1)/2}.
\]
Let $Q_{\mathrm{abs}}=1+a_3t^3+a_4t^4+a_5t^5+\frac12a_3^2t^6+a_3a_4t^7+\frac16a_3^3t^9$ and
$P_{\mathrm{abs}}=C_1t+\frac12(C_2+C_1^2)t^2+\frac16C_3t^3+\frac1{24}C_4t^4$. Then
\[
\begin{aligned}
E_N&=\mathcal I_{B'}[C_1t\,T_\omega]+\mathcal I_1[Q_{\mathrm{abs}}T_H]+e^{-w_0^2/4}\,\mathcal I_{1/2}[Q_{\mathrm{abs}}P_{\mathrm{abs}}]+2\mathfrak m,\\
E_D&=\mathcal I_{B'}[T_\omega]+e^{-w_0^2/4}\,\mathcal I_{1/2}[Q_{\mathrm{abs}}]+\mathfrak m,
\end{aligned}
\]
where $T_\omega,T_H$ are as in Section~\ref{sec:main} with $(\kappa_r,c_r)\to(A_r,C_r)$, and
\[
  \mathfrak m:=2\pi\,e^{-c_\ast/\theta}\Bigl(\frac{b}{2\pi\theta^3}\Bigr)^{1/2}
\]
is the minor-arc term $2\pi M_\chi/\sqrt{2\pi/B}$ with Corollary~\ref{cor:cstar}. (Since $\mathcal I_A$ integrates against $e^{-Aw^2/2}$, the central-arc terms carry the factor $B'^{-1/2}$ and the tail terms the
factor $\sqrt2$ automatically; the code evaluates $\mathcal I_A[P]=A^{-1/2}\,\E P(|Z_A|)$ with $Z_A\sim N(0,1/A)$.) With $N_0$, $M_2$ and
$\mathrm{Num}$ in scaled form,
\[
  \Delta=\frac{E_N+|\tilde N_0/D_0|E_D}{N_0-E_D}+\frac{|\mathrm{Num}|}{N_0},\qquad \varepsilon_k=-\log(1-\Delta/M_2).
\]
With the orders of magnitude noted after Lemma~\ref{lem:gap}, $E_N=O(\theta^{7/2})$, $E_D=O(\theta^2)$, $\tilde N_0/D_0=O(\theta^2)$ and
$\mathrm{Num}=O(\theta^4)$ (the tails and minor arcs being exponentially small), so $\varepsilon_k=O(\theta^{7/2})$, which is small compared
with $\theta q\asymp\theta^2$ for $\nu'$ in a bounded range. Theorem~\ref{thm:A} quantifies this.

%======================================================================
\section{The monotone majorant for \texorpdfstring{$\theta\le0.002$}{theta <= 0.002}}\label{app:regionII}
%======================================================================

Let $\theta_1=0.002$, $t=\sqrt\theta\le t_1=\sqrt{\theta_1}$, and fix a cell $[v_0,v_0']$ of $\nu'$. By~\eqref{eq:AC} and
$q=\theta e^{-\nu'}=t^2e^{-\nu'}$:
\begin{itemize}
\item $A_r=s_rb^{-r/2}t^{r-2}$, where $s_r\le r!\zeta_2+\theta_1\TV(f_r)$ and $b\ge F_2(\nu_{\min}-\theta_1)-\theta_1\TV(f_2)$;
\item $C_r=f_r(\nu)b^{-r/2}t^r$ with $f_r(\nu)=\nu^re^{-\nu}P_{r-1}(e^{-\nu})/(1-e^{-\nu})^r\le\nu^r\,q\,P_{r-1}(y)/(1-y)^r$,
  $y=e^{-\nu_{\min}}$ (Eulerian polynomials $P_{r-1}$ have nonnegative coefficients);
\item $C_1=c_1(\theta^3/b)^{1/2}$ with $c_1=1+\nu e^{-\nu'}/(1-q)$.
\end{itemize}
Each term of $E_N$, $E_D$, $\mathrm{Num}$ and $M_2-1$ is a polynomial with nonnegative coefficients in these quantities, except for the
explicitly signed terms of $M_2$ and $N_0$, which are bounded below termwise (for instance $N_0\ge1-\frac5{24}A_3^2$). Expanding, each term
is a nonnegative coefficient times $t^p\,q^{b_q}\,\nu^j$. After division by $\theta q=t^2q$ the monomial becomes
$t^P\nu^je^{-(b_q-1)\nu'}$, and at fixed $\nu'$, $\nu=\nu'+2\log(1/t)$. By the identity in Section~\ref{sec:proofA}, $t\mapsto t^P\nu^j$ is
nondecreasing on $(0,t_1]$ when $P>0$ and $\nu_{\min}>2j/P$ (for $j<0$ it is increasing whenever $P>0$). The program asserts these two
conditions for every non-constant monomial, evaluates each at $t=t_1$ with $\nu$ at the appropriate end ($\nu\le v_0'+\log(1/\theta_1)$ if $j\ge0$,
$\nu\ge\nu_{\min}$ if $j<0$), and bounds $e^{-(b_q-1)\nu'}$ at the appropriate end of the cell. The non-polynomial factors
$e^{-w_0^2/4}=e^{-0.04b/\theta}$ and $e^{-c_\ast/\theta}$ are handled by
$\frac{d}{d\theta}\log(\theta^{-p}e^{-a/\theta})=(a-p\theta)/\theta^2$, which is positive for $\theta<a/p$. The same method, applied to
the increments of Lemma~\ref{lem:B1} at the two points $k$ and $k+1$, gives Region II of Theorem~\ref{thm:B}.

%======================================================================
\begin{thebibliography}{99}

% VERIFY: volume/pages of the 1942 paper (J. Indian Math. Soc. (N.S.) 6 (1942)); title as cited in later literature.
\bibitem{ACG42}
F.~C.~Auluck, S.~Chowla and H.~Gupta, On the maximum value of the number of partitions of $n$ into $k$ parts,
\emph{J. Indian Math. Soc. (N.S.)} \textbf{6} (1942), 105--112.

\bibitem{Andrews}
G.~E.~Andrews, \emph{The Theory of Partitions}, Encyclopedia of Mathematics and its Applications, Vol.~2, Addison-Wesley, Reading, MA, 1976.

\bibitem{Can97}
E.~R.~Canfield, From recursions to asymptotics: on Szekeres' formula for the number of partitions,
\emph{Electron. J. Combin.} \textbf{4}(2) (1997), \#R6.

\bibitem{Erdos46}
P.~Erd\H{o}s, On some asymptotic formulas in the theory of partitions, \emph{Bull. Amer. Math. Soc.} \textbf{52} (1946), 185--188.

% VERIFY: volume and pages.
\bibitem{EL41}
P.~Erd\H{o}s and J.~Lehner, The distribution of the number of summands in the partitions of a positive integer,
\emph{Duke Math. J.} \textbf{8} (1941), 335--345.

\bibitem{Arb}
F.~Johansson, Arb: efficient arbitrary-precision midpoint-radius interval arithmetic,
\emph{IEEE Trans. Comput.} \textbf{66} (2017), 1281--1292. % VERIFY: page range

% VERIFY: title and pages.
\bibitem{Sz51}
G.~Szekeres, An asymptotic formula in the theory of partitions, \emph{Quart. J. Math. Oxford Ser. (2)} \textbf{2} (1951), 85--108.

% VERIFY: title and pages.
\bibitem{Sz53}
G.~Szekeres, Some asymptotic formulae in the theory of partitions.~II, \emph{Quart. J. Math. Oxford Ser. (2)} \textbf{4} (1953), 96--111.

% VERIFY: full bibliographic details of the 1990 paper.
\bibitem{Sz90}
G.~Szekeres, Asymptotic distribution of partitions by number and size of parts, in: \emph{Number Theory, Vol.~I (Budapest, 1987)},
Colloq. Math. Soc. J\'anos Bolyai \textbf{51}, North-Holland, Amsterdam, 1990, pp.~527--538.

\end{thebibliography}

\end{document}
